\documentclass[10pt]{article}
\usepackage[a4paper,margin=1in]{geometry}
\usepackage{amsmath,amssymb,amsfonts}
\usepackage{amsthm}
\usepackage{mathtools}
\usepackage{mathrsfs}
\usepackage{bm}
\usepackage{booktabs}
\usepackage{array}
\usepackage{multirow}
\usepackage{enumitem}
\usepackage{xcolor}
\usepackage{microtype}
\usepackage[colorlinks=true,linkcolor=blue,citecolor=blue,urlcolor=blue]{hyperref}

\theoremstyle{plain}
\newtheorem{theorem}{Theorem}[section]
\newtheorem{proposition}[theorem]{Proposition}
\newtheorem{lemma}[theorem]{Lemma}
\newtheorem{corollary}[theorem]{Corollary}
\theoremstyle{definition}
\newtheorem{definition}[theorem]{Definition}
\newtheorem{assumption}[theorem]{Assumption}
\theoremstyle{remark}
\newtheorem{remark}[theorem]{Remark}
\newtheorem{example}[theorem]{Example}

\newcommand{\R}{\mathbb{R}}
\newcommand{\N}{\mathbb{N}}
\newcommand{\Id}{\mathrm{Id}}
\newcommand{\Fix}{\operatorname{Fix}}
\newcommand{\Crit}{\operatorname{Crit}}
\newcommand{\spec}{\operatorname{spec}}
\newcommand{\dist}{\operatorname{dist}}
\newcommand{\diam}{\operatorname{diam}}
\newcommand{\rank}{\operatorname{rank}}
\newcommand{\cC}{\mathcal{C}}
\newcommand{\cB}{\mathcal{B}}
\newcommand{\cA}{\mathcal{A}}
\newcommand{\Om}{\Omega}
\newcommand{\Sig}{\Sigma}
\newcommand{\eps}{\varepsilon}
\newcommand{\grad}{\nabla}
\newcommand{\Ws}{W^{s}}
\newcommand{\Wu}{W^{u}}
\newcommand{\Wlocs}{W_{\mathrm{loc}}^{s}}
\newcommand{\Wlocu}{W_{\mathrm{loc}}^{u}}
\newcommand{\spr}{\rho}
\numberwithin{equation}{section}

\title{\textbf{Separation Laws for Flip Bifurcations and Endogenous
Switching Boundaries in Gradient Maps}}
\author{Zisong Yuan, Xiao Wen$^{\dagger}$\\\\
Department of Mathematics, Beihang University, Beijing, China\\\\
\normalsize E-mail: alexsongyzs@gmail.com (Z. Yuan), wenxiao@buaa.edu.cn (X. Wen)\\\\
$^{\dagger}$Corresponding author.}
\date{}

\begin{document}

\maketitle

\begin{abstract}
Gradient clipping is widely used to control large updates in
optimization, while finite-step gradient descent develops period-two
oscillations near the edge of linear stability.  Recent work has
analysed these oscillations via smooth flip bifurcation theory and
higher-order stability criteria, whereas piecewise-smooth bifurcation
theory describes generic interactions between period-doubling and
switching boundaries.  It remains unclear how a smooth
gradient-descent oscillation first interacts with the switching
surface generated endogenously by radial gradient clipping, and what
dynamical structure emerges after this contact.

For the finite-dimensional clipped gradient map
$F_{\eta,c}(x)=x-\eta\,\mathcal C_c^\diamond(\nabla V(x))$,
we analyse this transition near a nondegenerate strict local minimum
with a simple maximal Hessian eigenvalue.  Clipping is locally
inactive near the minimum, so the fixed-point stability threshold and
the local flip bifurcation coincide with those of ordinary gradient
descent.  Under a supercritical flip condition, the stable period-two
branch reaches the clipping boundary at
$\eta_{\mathrm{sc}}(c)=\eta_f+K_\diamond c^2+O(c^4)$ with
$\eta_f=2/\lambda_*$, where
$K_\diamond=\Gamma_*/(3\lambda_*^4\|v_*\|_\diamond^2)$
is determined explicitly by derivatives of the potential up to fourth
order and by the clipping norm.

Beyond the separation law, the gradient structure yields rigidity
absent from generic switching maps: the two points of every smooth
period-two orbit have opposite gradients and hence reach a radial
clipping boundary simultaneously; mixed smooth--clipped period-two
itineraries are impossible; and every fully clipped period-two orbit
has an exact unit Floquet multiplier and is therefore nonhyperbolic.
A solvable planar model realizes the full transition from a stable
smooth two-cycle to a neutral family of fully clipped period-two
cycles.
\end{abstract}

\noindent\textbf{Keywords:} gradient maps; gradient clipping; flip
bifurcation; period-doubling bifurcation; piecewise-smooth dynamical
systems; switching hypersurfaces; border collision; periodic orbits.

\vspace{1em}

% =========================================================
\section{Introduction}
\label{sec:introduction}
% =========================================================

Gradient descent with a finite step size is naturally a discrete
dynamical system,
\begin{equation}
x_{n+1}
=
x_n-\eta\nabla V(x_n),
\label{eq:intro-GD}
\end{equation}
and its behavior near the boundary of linear stability can differ
substantially from that suggested by the continuous gradient-flow
approximation; see, for example,
\cite{RobbinsMonro1951,Polyak1987,Nesterov2018,BottouCurtisNocedal2018}
for the classical optimization background.  In particular, when a nondegenerate local minimum
$x_*$ has maximal Hessian eigenvalue
\[
\lambda_*=
\lambda_{\max}(D^2V(x_*)),
\]
the classical local stability threshold is
\begin{equation}
\eta_f
=
\frac{2}{\lambda_*}.
\label{eq:intro-flip-threshold}
\end{equation}
At this value, the multiplier corresponding to the sharpest
eigendirection reaches $-1$, and, under the usual nondegeneracy
conditions, a flip bifurcation may occur.
This is the classical simple flip-bifurcation mechanism for smooth
maps; see, for example,
\cite{Kuznetsov2004,GuckenheimerHolmes1983,HaleKocak1991}.
The dynamics of gradient descent near this threshold have recently
received considerable attention, partly motivated by the
edge-of-stability phenomenon observed in large-step
neural-network optimization
\cite{CohenEtAl2021}.  Beyond
linear stability, higher-order derivatives of the objective determine
whether the loss of stability gives rise to stable or unstable
oscillations.  In the multivariate setting, the nonlinear coefficient
governing the local period-two branch depends not only on derivatives
along the maximal-curvature eigendirection, but also on its coupling to
the remaining Hessian eigendirections; see
\cite{MulayoffStich2026,Gan2026,Litman2026}.
These developments provide a
precise smooth bifurcation description of oscillatory gradient-descent
dynamics near a minimum.

Gradient clipping is a standard mechanism for controlling large
updates in optimization and has been studied from the viewpoints of
training stability, adaptivity, and convergence
\cite{PascanuMikolovBengio2013,ZhangEtAl2020,KoloskovaHendrikxStich2023,Goodfellow2016}.
The situation changes qualitatively when the gradient is clipped.
We consider the radial clipped gradient map
\begin{equation}
F_{\eta,c}(x)
=
x-\eta\mathcal C_c^\diamond(\nabla V(x)),
\label{eq:intro-clipped-map}
\end{equation}
where
\begin{equation}
\mathcal C_c^\diamond(g)
=
\min
\left\{
1,\frac{c}{\|g\|_\diamond}
\right\}g,
\qquad
c>0.
\label{eq:intro-clipping}
\end{equation}
Here $\|\cdot\|_\diamond$ is a smooth norm away from the origin.
The standard Euclidean clipping rule is recovered by taking
$\|\cdot\|_\diamond=\|\cdot\|_2$.

The clipping rule introduces the state-dependent switching
hypersurface
\begin{equation}
\Sigma_c
=
\left\{
x:
\|\nabla V(x)\|_\diamond=c
\right\}.
\label{eq:intro-switching-surface}
\end{equation}
Inside
\[
\Omega_c^-
=
\{x:\|\nabla V(x)\|_\diamond<c\},
\]
the map \eqref{eq:intro-clipped-map} coincides exactly with ordinary
gradient descent, whereas in
\[
\Omega_c^+
=
\{x:\|\nabla V(x)\|_\diamond>c\}
\]
the step magnitude is saturated.  Thus clipping converts a smooth
gradient map into a continuous piecewise-smooth dynamical system whose
switching boundary is not imposed externally: it is generated
endogenously by the same potential $V$ that generates the smooth
dynamics.  Piecewise-smooth maps and their border-collision
bifurcations have been studied extensively; see
\cite{Filippov1988,NusseYorke1992,diBernardo2008,diBernardoEtAl2008,
diBernardoHogan2010,Simpson2010,Simpson2016,LeineNijmeijer2004}.

This observation leads to a basic separation of mechanisms.  Since
\[
\nabla V(x_*)=0,
\]
for every fixed $c>0$ there exists a neighborhood of $x_*$ entirely
contained in $\Omega_c^-$.  Consequently, clipping cannot change the
local germ of the map at the minimum.  In particular, the fixed point,
its linearization, its first loss-of-stability threshold, the
center-manifold reduction, and the smooth flip coefficient are all
identical to those of ordinary gradient descent.  The first genuinely
clipping-induced effect must therefore occur at finite amplitude,
after the bifurcating orbit has grown far enough to reach
$\Sigma_c$.

The aim of this paper is to characterize this first
smooth-to-switching transition.

% ---------------------------------------------------------
\subsection{Main question}
\label{subsec:intro-main-question}
% ---------------------------------------------------------

Assume that $x_*$ is a nondegenerate strict local minimum,
that the maximal Hessian eigenvalue $\lambda_*$ is simple, and that
the ordinary gradient map undergoes a supercritical flip bifurcation
at
\[
\eta_f=\frac{2}{\lambda_*}.
\]
For
\[
\mu=\eta-\eta_f>0
\]
small, the smooth map possesses a stable period-two cycle
\begin{equation}
\mathcal O_2(\mu)
=
\{x_-(\mu),x_+(\mu)\}.
\label{eq:intro-cycle}
\end{equation}

The central question is:

\begin{quote}
At what parameter value does the smooth period-two branch first reach
the clipping boundary, and what dynamical structure is created when
this contact occurs?
\end{quote}

For a generic flip bifurcation, the oscillation amplitude grows on the
square-root scale
\[
\|x_\pm(\mu)-x_*\|
\asymp
\sqrt{\mu},
\]
whereas the local switching boundary generated by
\[
\|\nabla V(x)\|_\diamond=c
\]
lies at distance $O(c)$ from the critical point.  This suggests the
matching law
\[
\mu_{\mathrm{sc}}
\asymp
c^2.
\]
However, the main issue is not merely the exponent.  General
period-doubling--border-collision interactions are already known to
produce tangencies between smooth bifurcation and switching loci.
What is special here is the gradient-generated nature of both the
smooth map and the switching surface.

This additional structure allows us to derive the separation
coefficient explicitly from the local derivatives of $V$, to prove
that the two points of the period-two orbit reach the switching
surface simultaneously, and to determine a universal spectral
degeneracy in the fully clipped regime.

% ---------------------------------------------------------
\subsection{Main results}
\label{subsec:intro-results}
% ---------------------------------------------------------

Let
\[
H_*
=
D^2V(x_*),
\qquad
T_*
=
D^3V(x_*),
\qquad
Q_*
=
D^4V(x_*),
\]
and let
\[
H_*v_*=\lambda_*v_*,
\qquad
\|v_*\|_2=1.
\]
Define
\begin{equation}
g_*
=
T_*[v_*,v_*,\cdot]^\sharp
\label{eq:intro-g}
\end{equation}
and
\begin{equation}
\Gamma_*
=
3T_*[v_*,v_*,H_*^{-1}g_*]
-
Q_*[v_*,v_*,v_*,v_*].
\label{eq:intro-Gamma}
\end{equation}
The condition
\[
\Gamma_*>0
\]
is the supercritical flip condition used throughout the paper.

Our first result is a local identity principle.  For every fixed
$c>0$, the clipped and unclipped maps coincide on a neighborhood of
$x_*$.  Hence
\begin{equation}
DF_{\eta,c}(x_*)
=
I-\eta H_*,
\label{eq:intro-local-linearization}
\end{equation}
and the first local stability threshold is exactly
\begin{equation}
\boxed{
\eta_f
=
\frac{2}{\lambda_*},
}
\label{eq:intro-threshold-box}
\end{equation}
independently of $c$.  More strongly, the finite jets of the two maps
at $x_*$ agree, so the local center dynamics and smooth flip
coefficient are also clipping-independent.

Second, we derive the leading amplitude of the smooth period-two
branch.  For
\[
\mu=\eta-\eta_f\to0^+,
\]
one has
\begin{equation}
x_\pm(\mu)
=
x_*
\pm
A_*\sqrt{\mu}\,v_*
+
O(\mu),
\label{eq:intro-amplitude}
\end{equation}
where
\begin{equation}
\boxed{
A_*^2
=
\frac{3\lambda_*^2}{\Gamma_*}.
}
\label{eq:intro-A}
\end{equation}

Third, the gradient structure gives an exact identity that is stronger
than the local normal-form expansion.  For every genuine period-two
cycle of ordinary gradient descent,
\begin{equation}
\boxed{
\nabla V(x_+)
=
-\nabla V(x_-).
}
\label{eq:intro-gradient-balance}
\end{equation}
Thus the two points have exactly the same clipping exposure for every
norm:
\[
\|\nabla V(x_+)\|_\diamond
=
\|\nabla V(x_-)\|_\diamond.
\]
Consequently, the first interaction with a radial clipping boundary is
necessarily a simultaneous two-point contact.

Along the bifurcating branch, the common exposure satisfies
\begin{equation}
\|\nabla V(x_\pm(\mu))\|_\diamond
=
B_*\sqrt{\mu}
+
O(\mu^{3/2}),
\label{eq:intro-exposure}
\end{equation}
where
\begin{equation}
B_*
=
\lambda_*A_*
\|v_*\|_\diamond.
\label{eq:intro-B}
\end{equation}
Inverting the contact equation
\[
\|\nabla V(x_\pm(\mu_{\mathrm{sc}}))\|_\diamond=c
\]
gives our principal separation theorem:
\begin{equation}
\boxed{
\eta_{\mathrm{sc}}(c)-\eta_f
=
K_\diamond c^2
+
O(c^4),
\qquad
c\to0^+,
}
\label{eq:intro-main-law}
\end{equation}
with the explicit coefficient
\begin{equation}
\boxed{
K_\diamond
=
\frac{\Gamma_*}{
3\lambda_*^4
\|v_*\|_\diamond^2
}.
}
\label{eq:intro-K}
\end{equation}

Equivalently, if
\[
\{v_i\}_{i=1}^d
\]
is an orthonormal Hessian eigenbasis with eigenvalues
$\lambda_i$, then
\begin{equation}
K_\diamond
=
\frac{
3\displaystyle\sum_{i=1}^d
\frac{
T_*[v_*,v_*,v_i]^2
}{
\lambda_i
}
-
Q_*[v_*^4]
}{
3\lambda_*^4
\|v_*\|_\diamond^2
}.
\label{eq:intro-K-eigen}
\end{equation}

Thus the same high-order local quantity that controls the birth and
stability of the smooth oscillation also determines the leading
distance in parameter space to the first clipping event.

The radial structure imposes further restrictions after contact.
A period-two orbit of the clipped map cannot have one point in
$\Omega_c^-$ and the other in $\Omega_c^+$.  Hence a mixed
smooth--clipped period-two itinerary is impossible.  If a period-two
orbit persists beyond contact, both points must lie in the clipped
region and satisfy the exact amplitude constraint
\begin{equation}
\boxed{
\|y_+-y_-\|_\diamond
=
\eta c.
}
\label{eq:intro-clipped-amplitude}
\end{equation}

Finally, every fully clipped period-two orbit is necessarily
nonhyperbolic.  If
\[
M_c
=
DF_{\eta,c}(y_-)
DF_{\eta,c}(y_+)
\]
denotes its monodromy matrix, then
\begin{equation}
\boxed{
1\in\sigma(M_c).
}
\label{eq:intro-unit-multiplier}
\end{equation}
This unit multiplier follows from the fact that radial clipping
annihilates infinitesimal changes in gradient magnitude while
retaining only tangent variations of the gradient direction.

A completely solvable planar gradient model shows that this
degeneracy can have a genuine nonlinear manifestation: after the
simultaneous switching contact, the isolated attracting smooth
period-two cycle is replaced locally by a one-parameter family of
fully clipped period-two cycles, transversely attracting but neutral
along the family.

\subsection{Relation to existing bifurcation theory}

The present analysis lies at the intersection of two established
lines of work.

The first is smooth local bifurcation theory for maps.
Near a simple multiplier $-1$, center-manifold and normal-form
reductions yield the classical flip bifurcation and the
square-root scaling of the emerging period-two branch; see,
for example,
\cite{Kuznetsov2004,GuckenheimerHolmes1983,ChowHale1982,
Wiggins2003,Arnold1983,Sotomayor1973,HaleKocak1991,PalisDeMelo1982}.

More recently, nonlinear oscillatory dynamics of gradient descent
near the edge of stability have been studied directly.
Mulayoff and Stich \cite{MulayoffStich2026} derived a multivariate
stability criterion for nonlinear gradient-descent oscillations
near a minimum in terms of higher-order derivatives of the
objective.
Gan \cite{Gan2026} formulated edge-of-stability dynamics in a
bifurcation-theoretic framework, in which the relevant normal
dynamics undergo a flip bifurcation governed by the corresponding
first Lyapunov coefficient.
Relatedly, Litman \cite{Litman2026} characterized fixed points and
period-two orbits of gradient descent through an exact two-step
variational formulation and obtained a local reduction in terms
of the half-amplitude of the oscillation.

Some ingredients of the smooth analysis developed below therefore
belong to standard flip-bifurcation theory or have close analogues
in recent analyses of nonlinear gradient-descent dynamics.
We nevertheless include self-contained derivations, both to fix
notation and because the resulting coefficients enter explicitly
into the subsequent switching analysis.

The second relevant line of work is the theory of piecewise-smooth
maps and border-collision bifurcations
\cite{NusseYorke1995,BanerjeeGrebogi1999,Feigin1970,diBernardo2008,
Simpson2010,Simpson2016}.
In particular, Simpson and Meiss
\cite{SimpsonMeiss2008,SimpsonMeiss2009,SimpsonMeiss2009Nonlinearity,
SimpsonMeiss2010,SimpsonMeiss2012} studied the codimension-two
interaction between period-doubling and border-collision
bifurcations in general piecewise-smooth continuous maps. They
showed that the border-collision locus of the period-doubled
solution can be tangent to the period-doubling locus near the
codimension-two point. Thus quadratic separation between smooth
bifurcation and switching loci is already a feature of generic
period-doubling--border-collision interactions.

Our contribution is not the generic existence of such a tangency,
nor the smooth flip mechanism itself. Rather, we exploit the
additional variational structure of clipped gradient maps.
The switching surface
\[
\Sigma_c=\{x:\|\nabla V(x)\|_\diamond=c\}
\]
is generated endogenously by the same potential that determines
the smooth gradient dynamics. Moreover, the period-two equations
impose exact gradient-balance identities.

These additional constraints yield structural consequences that
are absent from a generic externally switched map. In particular,
we obtain:
\begin{enumerate}
    \item clipping-independent local flip dynamics;
    \item an explicit tensor formula for the flip-to-switching
    separation coefficient;
    \item simultaneous switching contact of the two points of a
    smooth period-two orbit;
    \item nonexistence of mixed smooth--clipped period-two
    itineraries;
    \item an exact unit Floquet multiplier for every fully clipped
    period-two orbit.
\end{enumerate}

It is these gradient-induced rigidity properties, rather than the
standard smooth flip bifurcation or the generic tangency mechanism,
that form the main focus of the present work.
% ---------------------------------------------------------
\subsection{Organization of the paper}
\label{subsec:intro-organization}
% ---------------------------------------------------------

The paper is organized as follows.

Section~\ref{sec:setting} introduces the clipped gradient map,
the switching hypersurface, and the standing assumptions.

Section~\ref{sec:local-structure} proves the local identity between
clipped and ordinary gradient descent and establishes the
clipping-independent stability threshold.

Section~\ref{sec:flip-bifurcation} derives the supercritical
period-two branch and its explicit leading amplitude.

Section~\ref{sec:switching-geometry} develops the exact gradient
balance and switching geometry along the period-two branch.

Section~\ref{sec:separation-law} proves the quadratic
flip-to-switching separation law and derives its explicit coefficient.

Section~\ref{sec:post-contact} studies the fully clipped regime and
proves the universal unit-multiplier theorem.

Section~\ref{sec:planar-model} presents a completely solvable planar
model that realizes the full transition analytically.


% =========================================================
\section{Problem setting and standing hypotheses}
\label{sec:setting}
% =========================================================

Throughout the paper, $\langle\cdot,\cdot\rangle$ and
$\|\cdot\|_2$ denote the standard Euclidean inner product and norm
on $\mathbb{R}^d$, respectively. All gradients, Hessians, and spectral
quantities are understood with respect to this Euclidean structure.

Let $\mathcal U\subset\mathbb{R}^d$ be an open set and let
\[
V:\mathcal U\to\mathbb{R}
\]
be a smooth potential. For a step size $\eta>0$, the associated
gradient-descent map is
\begin{equation}
G_\eta(x)
=
x-\eta\nabla V(x).
\label{eq:smooth-gd-map}
\end{equation}

Our aim is to study how a smooth flip bifurcation of
\eqref{eq:smooth-gd-map} interacts with the switching boundary
generated intrinsically by gradient clipping.

% ---------------------------------------------------------
\subsection{Clipped gradient maps}
\label{subsec:clipped-map}
% ---------------------------------------------------------

Let $\|\cdot\|_\diamond$ be a norm on $\mathbb{R}^d$.
For a clipping threshold $c>0$, define
\begin{equation}
\mathcal C_c^\diamond(g)
=
\begin{cases}
g,
&
\|g\|_\diamond\leq c,
\\[1mm]
\dfrac{c}{\|g\|_\diamond}\,g,
&
\|g\|_\diamond>c,
\end{cases}
\qquad
g\in\mathbb{R}^d.
\label{eq:clipping-operator}
\end{equation}
Equivalently,
\[
\mathcal C_c^\diamond(g)
=
\min\left\{
1,\frac{c}{\|g\|_\diamond}
\right\}g,
\]
with the convention $\mathcal C_c^\diamond(0)=0$.

The clipped gradient map is
\begin{equation}
F_{\eta,c}(x)
=
x-\eta\mathcal C_c^\diamond(\nabla V(x)).
\label{eq:clipped-gradient-map}
\end{equation}

The standard Euclidean gradient-clipping rule corresponds to
\[
\|\cdot\|_\diamond=\|\cdot\|_2.
\]

Associated with \eqref{eq:clipped-gradient-map} is the endogenous
switching hypersurface
\begin{equation}
\Sigma_c
=
\left\{
x\in\mathcal U:
\|\nabla V(x)\|_\diamond=c
\right\}.
\label{eq:switching-hypersurface}
\end{equation}
It separates the local phase space into the smooth region
\begin{equation}
\Omega_c^-
=
\left\{
x\in\mathcal U:
\|\nabla V(x)\|_\diamond<c
\right\}
\label{eq:smooth-region}
\end{equation}
and the clipped region
\begin{equation}
\Omega_c^+
=
\left\{
x\in\mathcal U:
\|\nabla V(x)\|_\diamond>c
\right\}.
\label{eq:clipped-region}
\end{equation}

Accordingly,
\begin{equation}
F_{\eta,c}(x)
=
G_\eta(x)
=
x-\eta\nabla V(x),
\qquad
x\in\Omega_c^-,
\label{eq:smooth-branch}
\end{equation}
whereas in $\Omega_c^+$,
\begin{equation}
F_{\eta,c}(x)
=
x
-
\eta c
\frac{\nabla V(x)}
{\|\nabla V(x)\|_\diamond}.
\label{eq:clipped-branch}
\end{equation}

Thus $F_{\eta,c}$ is continuous and piecewise smooth, while its
switching geometry is not prescribed independently of the dynamics:
the boundary $\Sigma_c$ is generated by the same potential $V$
that defines the gradient map.

% ---------------------------------------------------------
\subsection{The distinguished critical point}
\label{subsec:critical-point}
% ---------------------------------------------------------

We focus on the dynamics near an isolated nondegenerate local
minimum $x_*$ of $V$.

Set
\begin{equation}
H_*:=D^2V(x_*).
\label{eq:H-star}
\end{equation}
Since $H_*$ is symmetric, choose an orthonormal eigenbasis
$\{v_i\}_{i=1}^d$ satisfying
\begin{equation}
H_*v_i=\lambda_i v_i,
\qquad
0<\lambda_1\leq\lambda_2\leq\cdots\leq\lambda_d.
\label{eq:H-eigendecomposition}
\end{equation}

We denote the maximal eigenvalue and its corresponding unit
eigenvector by
\begin{equation}
\lambda_*:=\lambda_d,
\qquad
v_*:=v_d.
\label{eq:critical-eigenpair}
\end{equation}

The linearization of the smooth gradient map at $x_*$ is
\begin{equation}
DG_\eta(x_*)
=
I-\eta H_*,
\label{eq:GD-linearization}
\end{equation}
and hence its multipliers are
\begin{equation}
\rho_i(\eta)
=
1-\eta\lambda_i,
\qquad
i=1,\ldots,d.
\label{eq:GD-multipliers}
\end{equation}

The distinguished step size at which the maximal-curvature
multiplier reaches $-1$ is
\begin{equation}
\eta_f
:=
\frac{2}{\lambda_*}.
\label{eq:eta-f}
\end{equation}

At $\eta=\eta_f$,
\begin{equation}
\rho_d(\eta_f)=-1.
\label{eq:critical-multiplier}
\end{equation}

The parameter
\begin{equation}
\mu:=\eta-\eta_f
\label{eq:mu-definition}
\end{equation}
will be used as the local bifurcation parameter throughout the paper.

% ---------------------------------------------------------
\subsection{Higher-order tensors and the flip coefficient}
\label{subsec:flip-coefficient-notation}
% ---------------------------------------------------------

We write
\begin{equation}
T_*:=D^3V(x_*),
\qquad
Q_*:=D^4V(x_*),
\label{eq:TQ-definitions}
\end{equation}
viewed as symmetric multilinear forms.

For $u,w\in\mathbb{R}^d$, let
\begin{equation}
T_*[u,w,\cdot]^\sharp
\in\mathbb{R}^d
\label{eq:sharp-definition}
\end{equation}
denote the Euclidean Riesz representative of the linear functional
$z\mapsto T_*[u,w,z]$; namely,
\begin{equation}
\left\langle
T_*[u,w,\cdot]^\sharp,z
\right\rangle
=
T_*[u,w,z]
\qquad
\text{for all }z\in\mathbb{R}^d.
\label{eq:sharp-identity}
\end{equation}

Define
\begin{equation}
g_*
:=
T_*[v_*,v_*,\cdot]^\sharp
\label{eq:g-star}
\end{equation}
and
\begin{equation}
h_*
:=
H_*^{-1}g_*.
\label{eq:h-star}
\end{equation}

The nonlinear quantity governing the local flip bifurcation is
\begin{equation}
\Gamma_*
:=
3T_*[v_*,v_*,h_*]
-
Q_*[v_*,v_*,v_*,v_*].
\label{eq:Gamma-star}
\end{equation}

Equivalently, using the eigenbasis
\eqref{eq:H-eigendecomposition},
\begin{equation}
\Gamma_*
=
3
\sum_{i=1}^d
\frac{
T_*[v_*,v_*,v_i]^2
}{
\lambda_i
}
-
Q_*[v_*,v_*,v_*,v_*].
\label{eq:Gamma-eigenbasis}
\end{equation}

With the flip-normal-form convention used in this paper, define
\begin{equation}
\ell_*
:=
\frac{\eta_f}{6}\Gamma_*
=
\frac{\Gamma_*}{3\lambda_*}.
\label{eq:ell-star}
\end{equation}
Thus $\ell_*>0$ corresponds to the supercritical case in which a
stable small-amplitude period-two branch is created for
$\eta>\eta_f$.

The expression \eqref{eq:Gamma-eigenbasis} is the multivariate
gradient-map form of the nonlinear flip criterion: the cubic
coupling of the critical direction to all Hessian eigendirections
contributes to the sign of the bifurcation coefficient, rather than
only the derivatives taken purely along $v_*$.

% ---------------------------------------------------------
\subsection{Standing hypotheses}
\label{subsec:standing-hypotheses}
% ---------------------------------------------------------

The following assumptions are imposed unless stated otherwise.

\begin{enumerate}[
label=\textbf{(H\arabic*)},
leftmargin=2.8em,
itemsep=0.8em
]

\item
\label{H1}
\textbf{Regularity.}
The potential satisfies
\[
V\in C^7(\mathcal U,\mathbb{R}).
\]

\item
\label{H2}
\textbf{Nondegenerate strict local minimum.}
The distinguished point $x_*\in\mathcal U$ satisfies
\begin{equation}
\nabla V(x_*)=0,
\qquad
H_*=D^2V(x_*)\succ0.
\label{eq:H2}
\end{equation}

\item
\label{H3}
\textbf{Simple maximal curvature.}
The largest Hessian eigenvalue is simple:
\begin{equation}
\lambda_{d-1}<\lambda_d=\lambda_*.
\label{eq:H3}
\end{equation}

\item
\label{H4}
\textbf{Supercritical flip nondegeneracy.}
The coefficient defined in \eqref{eq:Gamma-star} satisfies
\begin{equation}
\Gamma_*>0,
\label{eq:H4}
\end{equation}
or equivalently
\[
\ell_*>0.
\]

\item
\label{H5}
\textbf{Smooth clipping geometry.}
The clipping norm $\|\cdot\|_\diamond$ is of class $C^3$ on
$\mathbb{R}^d\setminus\{0\}$.
\end{enumerate}

% ---------------------------------------------------------
\subsection{Immediate spectral consequences}
\label{subsec:spectral-consequences}
% ---------------------------------------------------------

Assumptions \textbf{(H2)--(H3)} imply that the spectral hypotheses
required for a simple codimension-one flip bifurcation are automatic
for the gradient map.

Indeed, for $i=1,\ldots,d-1$,
\[
0<\lambda_i<\lambda_*,
\]
and therefore, at $\eta=\eta_f$,
\begin{equation}
-1
<
1-\frac{2\lambda_i}{\lambda_*}
<
1.
\label{eq:noncritical-spectrum}
\end{equation}
Hence $-1$ is the unique multiplier of $DG_{\eta_f}(x_*)$ on the
unit circle.

Moreover,
\begin{equation}
\frac{d}{d\eta}\rho_d(\eta)
=
-\lambda_*,
\label{eq:eigenvalue-crossing}
\end{equation}
so that
\begin{equation}
\left.
\frac{d}{d\eta}\rho_d(\eta)
\right|_{\eta=\eta_f}
=
-\lambda_*
\neq0.
\label{eq:transversal-crossing}
\end{equation}
Thus the critical multiplier crosses $-1$ transversally as $\eta$
passes through $\eta_f$.

Thus, under the spectral assumptions \textbf{(H2)--(H3)}, the sign of
$\Gamma_*$ provides the remaining nonlinear nondegeneracy condition
that distinguishes the supercritical and subcritical flip cases.

% ---------------------------------------------------------
\subsection{Local form of the switching geometry}
\label{subsec:local-switching-preview}
% ---------------------------------------------------------

We record one elementary observation that motivates the later
switching analysis.

Let
\[
x=x_*+sv_*+r(s),
\qquad
r(s)=O(s^2),
\qquad
s\to0.
\]
Since $\nabla V(x_*)=0$ and $H_*v_*=\lambda_*v_*$,
Taylor expansion gives
\begin{equation}
\nabla V(x)
=
\lambda_*s\,v_*
+
O(s^2).
\label{eq:gradient-critical-expansion}
\end{equation}

By positive homogeneity and smoothness of the clipping norm away
from the origin,
\begin{equation}
\|\nabla V(x)\|_\diamond
=
\lambda_*
\|v_*\|_\diamond
|s|
+
O(s^2).
\label{eq:norm-critical-expansion}
\end{equation}

In particular, the leading-order displacement of the switching
boundary from $x_*$ in the critical direction is nondegenerate.

For the Euclidean clipping rule,
$\|v_*\|_2=1$, and
\eqref{eq:norm-critical-expansion} reduces to
\begin{equation}
\|\nabla V(x)\|_2
=
\lambda_*|s|
+
O(s^2).
\label{eq:euclidean-critical-expansion}
\end{equation}

 A precise expansion of $\Sigma_c$ along the center manifold will be
derived later in Section~\ref{sec:switching-geometry}.

% ---------------------------------------------------------
\subsection{Scope of the present paper}
\label{subsec:scope}
% ---------------------------------------------------------

The positive-definiteness assumption in \textbf{(H2)} deliberately
restricts the present work to an isolated strict local minimum.
This excludes the additional unit multiplier $+1$ that occurs when
the Hessian has a nontrivial kernel, as in Morse--Bott manifolds of
minimizers.

The purpose of this restriction is to isolate the interaction between
a simple smooth flip bifurcation and the endogenous switching
boundary generated by clipping.  Extensions involving manifolds of
minimizers, multiple critical multipliers, or degenerate flip
coefficients require a higher-dimensional center dynamics and are
outside the scope of the present paper.



% =========================================================
\section{Local structure and clipping-independent stability}
\label{sec:local-structure}
% =========================================================

We first establish a structural fact used throughout the paper: for
every fixed $c>0$, clipping is completely inactive in a sufficiently
small neighborhood of $x_*$, so the clipped and unclipped gradient
maps have identical local germs at $x_*$.  Consequently, clipping
cannot alter the linear stability threshold, the local invariant
structures, or any smooth local normal-form coefficient at $x_*$;
its influence can appear only when an orbit reaches a
finite-amplitude region intersecting the switching hypersurface
$\Sigma_c$.

% ---------------------------------------------------------
\subsection{Fixed points}
\label{subsec:fixed-points}
% ---------------------------------------------------------

We begin with a global observation.

\begin{proposition}[Fixed-point invariance]
\label{prop:fixed-point-invariance}
Let $\eta>0$ and $c>0$. Then
\begin{equation}
\operatorname{Fix}(F_{\eta,c})
=
\operatorname{Crit}(V)
:=
\left\{
x\in\mathcal U:
\nabla V(x)=0
\right\}.
\label{eq:fixed-critical-equivalence}
\end{equation}
In particular, the set of fixed points is independent of both
$\eta$ and $c$.
\end{proposition}

\begin{proof}
If $\nabla V(x)=0$, then
\[
\mathcal C_c^\diamond(\nabla V(x))=0,
\]
and hence
\[
F_{\eta,c}(x)
=
x-\eta\mathcal C_c^\diamond(\nabla V(x))
=
x.
\]
Thus every critical point of $V$ is a fixed point of $F_{\eta,c}$.

Conversely, suppose that
\[
F_{\eta,c}(x)=x.
\]
Since $\eta>0$, we obtain
\[
\mathcal C_c^\diamond(\nabla V(x))=0.
\]
By the definition of the radial clipping operator
\eqref{eq:clipping-operator},
\[
\mathcal C_c^\diamond(g)=0
\quad\Longleftrightarrow\quad
g=0.
\]
Therefore
\[
\nabla V(x)=0.
\]
This proves \eqref{eq:fixed-critical-equivalence}.
\end{proof}

\begin{remark}
\label{rem:no-spurious-fixed-points}
Proposition~\ref{prop:fixed-point-invariance} shows that norm clipping
does not create spurious equilibria.  All dynamical effects of
clipping therefore arise from the behavior of nonstationary
trajectories rather than from a modification of the critical-point
set of the potential.
\end{remark}

% ---------------------------------------------------------
\subsection{Quantitative inactivity of clipping}
\label{subsec:local-inactivity}
% ---------------------------------------------------------

The next lemma gives a quantitative local estimate.

For a linear map $A:\mathbb R^d\to\mathbb R^d$, define the mixed
operator norm
\begin{equation}
\|A\|_{2\to\diamond}
:=
\sup_{\|z\|_2=1}
\|Az\|_\diamond.
\label{eq:mixed-operator-norm}
\end{equation}

\begin{lemma}[Local gradient estimate]
\label{lem:local-gradient-estimate}
Under \textbf{(H1)--(H2)}, there exist $r_0>0$ and $L_*>0$ such that
\[
\overline{B_{r_0}(x_*)}
\subset\mathcal U
\]
and
\begin{equation}
\|\nabla V(x)\|_\diamond
\leq
L_*\|x-x_*\|_2
\qquad
\text{for all }
x\in B_{r_0}(x_*).
\label{eq:local-gradient-bound}
\end{equation}
\end{lemma}

\begin{proof}
Since $\mathcal U$ is open and $x_*\in\mathcal U$, choose
$r_0>0$ sufficiently small so that
\[
\overline{B_{r_0}(x_*)}
\subset\mathcal U.
\]
By \textbf{(H1)}, the Hessian $D^2V$ is continuous. Hence
\begin{equation}
L_*
:=
\sup_{y\in\overline{B_{r_0}(x_*)}}
\|D^2V(y)\|_{2\to\diamond}
<\infty.
\label{eq:L-star-definition}
\end{equation}

For $x\in B_{r_0}(x_*)$, the line segment joining $x_*$ to $x$
lies in $B_{r_0}(x_*)$. Since $\nabla V(x_*)=0$, the fundamental
theorem of calculus gives
\begin{align}
\nabla V(x)
&=
\nabla V(x)-\nabla V(x_*)
\nonumber\\
&=
\int_0^1
D^2V\bigl(x_*+t(x-x_*)\bigr)(x-x_*)\,dt.
\label{eq:gradient-integral}
\end{align}
Therefore
\begin{align}
\|\nabla V(x)\|_\diamond
&\leq
\int_0^1
\left\|
D^2V\bigl(x_*+t(x-x_*)\bigr)
\right\|_{2\to\diamond}
\|x-x_*\|_2\,dt
\nonumber\\
&\leq
L_*\|x-x_*\|_2.
\end{align}
This proves \eqref{eq:local-gradient-bound}.
\end{proof}

For every fixed $c>0$, define
\begin{equation}
r_c
:=
\min
\left\{
r_0,
\frac{c}{2L_*}
\right\}.
\label{eq:rc-definition}
\end{equation}

\begin{corollary}[Clipping inactivity]
\label{cor:clipping-inactivity}
For every $c>0$,
\begin{equation}
B_{r_c}(x_*)
\subset\Omega_c^-.
\label{eq:ball-inside-smooth-region}
\end{equation}
Equivalently,
\begin{equation}
B_{r_c}(x_*)\cap\Sigma_c
=
\varnothing.
\label{eq:no-switching-near-critical}
\end{equation}
\end{corollary}

\begin{proof}
If $x\in B_{r_c}(x_*)$, then by
Lemma~\ref{lem:local-gradient-estimate},
\[
\|\nabla V(x)\|_\diamond
\leq
L_*\|x-x_*\|_2
<
L_*r_c
\leq
\frac{c}{2}
<
c.
\]
Hence $x\in\Omega_c^-$.
\end{proof}

\begin{remark}[Dependence on the clipping threshold]
\label{rem:shrinking-neighborhood}
The neighborhood in Corollary~\ref{cor:clipping-inactivity} is not
uniform as $c\to0^+$.  The explicit choice
\eqref{eq:rc-definition} has
\[
r_c=O(c).
\]
This shrinking local scale is essential for the later
smooth-to-switching analysis: although clipping is invisible in a
neighborhood of $x_*$ for every fixed $c>0$, the boundary of that
neighborhood approaches the critical point as $c\to0^+$.
\end{remark}

% ---------------------------------------------------------
\subsection{Local identity of the clipped and smooth maps}
\label{subsec:local-identity}
% ---------------------------------------------------------

We can now state the basic structural theorem.

\begin{theorem}[Local identity theorem]
\label{thm:local-identity}
Assume \textbf{(H1)--(H2)}. For every fixed $c>0$, there exists
$r_c>0$ such that, for every $\eta>0$,
\begin{equation}
F_{\eta,c}(x)
=
G_\eta(x)
=
x-\eta\nabla V(x)
\qquad
\text{for all }
x\in B_{r_c}(x_*).
\label{eq:local-map-identity}
\end{equation}
The radius $r_c$ may be chosen independently of $\eta$.
\end{theorem}

\begin{proof}
By Corollary~\ref{cor:clipping-inactivity},
\[
B_{r_c}(x_*)
\subset\Omega_c^-.
\]
Hence, for every $x\in B_{r_c}(x_*)$,
\[
\|\nabla V(x)\|_\diamond<c.
\]
By definition of the clipping operator,
\[
\mathcal C_c^\diamond(\nabla V(x))
=
\nabla V(x).
\]
Substitution into \eqref{eq:clipped-gradient-map} yields
\[
F_{\eta,c}(x)
=
x-\eta\nabla V(x)
=
G_\eta(x).
\]
The construction of $r_c$ in \eqref{eq:rc-definition} does not
involve $\eta$, so the same neighborhood works for every $\eta>0$.
\end{proof}

Theorem~\ref{thm:local-identity} is stronger than equality of the
linearizations: the two maps agree as functions on an open
neighborhood of $x_*$.

% ---------------------------------------------------------
\subsection{Jet invariance and local normal forms}
\label{subsec:jet-invariance}
% ---------------------------------------------------------

The local identity immediately implies equality of all derivatives
allowed by the regularity of $V$.

\begin{corollary}[Local jet and normal-form invariance]
\label{cor:jet-invariance}
\label{cor:flip-coefficient-independence}
Under \textbf{(H1)--(H4)}, for every fixed $c>0$ and $\eta>0$:
\begin{enumerate}[label=\textnormal{(\roman*)}]
\item
every local Taylor coefficient of $F_{\eta,c}$ and $G_\eta$ at
$x_*$ that is defined under the assumed regularity is identical;
in particular,
\begin{equation}
DF_{\eta,c}(x_*)
=
I-\eta H_*;
\label{eq:clipped-linearization}
\end{equation}
\item
the local center-manifold equations and all smooth flip
normal-form coefficients of $F_{\eta,c}$ at $x_*$ are identical to
those of $G_\eta$; in particular, the coefficient $\ell_*$
defined in \eqref{eq:ell-star} is independent of the clipping
threshold $c$.
\end{enumerate}
\end{corollary}

\begin{proof}
The two maps coincide on a neighborhood of $x_*$ by
Theorem~\ref{thm:local-identity}, so all derivatives at $x_*$ that
are defined coincide.  Smooth normal-form coefficients depend only
on a finite jet at the bifurcation point; hence the center-manifold
invariance equations and reduced normal forms are identical.
\end{proof}

\begin{remark}
\label{rem:local-versus-finite-amplitude}
Corollary~\ref{cor:flip-coefficient-independence} is a central
structural distinction in the present problem.  Clipping cannot
modify the birth mechanism of a sufficiently small smooth
period-two orbit.  It can affect that orbit only after its amplitude
has grown sufficiently for one of its points to reach $\Sigma_c$.
The latter is a finite-amplitude event and will be analyzed in
Sections~\ref{sec:switching-geometry} and
\ref{sec:separation-law}.
\end{remark}

% ---------------------------------------------------------
\subsection{Exact local stability threshold}
\label{subsec:local-stability-threshold}
% ---------------------------------------------------------

We next determine the stability of $x_*$.  The threshold
$\eta_f=2/\lambda_*$ below is the classical maximal step size of
(unclipped) gradient descent \cite{Polyak1987,Nesterov2018};
Theorem~\ref{thm:local-stability-threshold} shows that it is
preserved exactly by clipping.

\begin{theorem}[Clipping-independent local stability threshold]
\label{thm:local-stability-threshold}
Assume \textbf{(H1)--(H3)} and let
\[
\eta_f=\frac{2}{\lambda_*}.
\]
Then, for every fixed $c>0$, the following statements hold.

\begin{enumerate}[label=\textnormal{(\roman*)}]

\item
If
\begin{equation}
0<\eta<\eta_f,
\label{eq:stable-eta-range}
\end{equation}
then $x_*$ is a locally asymptotically stable fixed point of
$F_{\eta,c}$.

\item
If
\begin{equation}
\eta>\eta_f,
\label{eq:unstable-eta-range}
\end{equation}
then $x_*$ is unstable for $F_{\eta,c}$.

\item
At
\begin{equation}
\eta=\eta_f,
\end{equation}
the fixed point is nonhyperbolic.  Its unique multiplier on the unit
circle is the simple multiplier $-1$ associated with $v_*$.

\end{enumerate}

Consequently, the first local stability threshold is exactly
\begin{equation}
\boxed{
\eta_f
=
\frac{2}{\lambda_{\max}(H_*)}
}
\label{eq:boxed-stability-threshold}
\end{equation}
and is independent of the clipping threshold $c$.
\end{theorem}

\begin{proof}
By Corollary~\ref{cor:jet-invariance},
\[
DF_{\eta,c}(x_*)
=
I-\eta H_*.
\]
Since $H_*$ is symmetric positive definite, the multipliers are
\[
\rho_i(\eta)
=
1-\eta\lambda_i,
\qquad
i=1,\ldots,d.
\]

Suppose first that
\[
0<\eta<\frac{2}{\lambda_*}.
\]
Since $0<\lambda_i\leq\lambda_*$,
\[
0<\eta\lambda_i<2,
\]
and hence
\[
-1
<
1-\eta\lambda_i
<
1
\qquad
\text{for every }i.
\]
Thus
\[
\max_i|\rho_i(\eta)|<1.
\]
The fixed point is therefore locally asymptotically stable by the
standard linearization theorem for discrete dynamical systems.

Now suppose that
\[
\eta>\frac{2}{\lambda_*}.
\]
Then
\[
\rho_d(\eta)
=
1-\eta\lambda_*
<
-1,
\]
so
\[
|\rho_d(\eta)|>1.
\]
Hence $x_*$ is unstable.

Finally, if
\[
\eta=\eta_f=\frac{2}{\lambda_*},
\]
then
\[
\rho_d(\eta_f)=-1.
\]
By \textbf{(H3)}, for every $i<d$,
\[
0<\lambda_i<\lambda_*,
\]
and therefore
\[
-1
<
1-\frac{2\lambda_i}{\lambda_*}
<
1.
\]
Thus $-1$ is the unique multiplier on the unit circle, and it is
simple.
\end{proof}

\begin{remark}
\label{rem:boundary-inconclusive-linearly}
Theorem~\ref{thm:local-stability-threshold} deliberately makes no
stability claim at $\eta=\eta_f$.  Linearization is inconclusive
when the critical multiplier has modulus one.  The nonlinear
behavior at this parameter is governed by the flip coefficient
$\ell_*$ and will be analyzed in Section~\ref{sec:flip-bifurcation}.
\end{remark}

% ---------------------------------------------------------
\subsection{Local invariant structures}
\label{subsec:local-invariant-structures}
% ---------------------------------------------------------

The local identity also transfers invariant structures from the
smooth gradient map to the clipped map.  For a map $T$ and a radius
$r>0$, define the local stable set relative to $B_r(x_*)$ by
\begin{equation}
W^s_r(x_*;T)
:=
\left\{
x\in B_r(x_*):
T^n(x)\in B_r(x_*)\ \forall n\geq0,
\quad
T^n(x)\to x_*
\right\}.
\label{eq:relative-local-stable-set}
\end{equation}

\begin{proposition}[Equality of local stable dynamics]
\label{prop:local-stable-dynamics}
Let $c>0$, and choose $0<r\leq r_c$.
Then, for every $\eta>0$,
\begin{equation}
W^s_r(x_*;F_{\eta,c})
=
W^s_r(x_*;G_\eta).
\label{eq:stable-set-equality}
\end{equation}
\end{proposition}

\begin{proof}
The two maps coincide on $B_r(x_*)$ by
Theorem~\ref{thm:local-identity}, so the defining conditions in
\eqref{eq:relative-local-stable-set} are equivalent for the two
maps.
\end{proof}

When $x_*$ is hyperbolic, this equality gives equality of the
corresponding stable-manifold germs.  If, in addition,
\begin{equation}
\det(I-\eta H_*)\neq0,
\label{eq:local-diffeomorphism-condition}
\end{equation}
then both maps are local diffeomorphisms near $x_*$, and the same
argument gives equality of their local unstable-manifold germs.

At the critical parameter $\eta=\eta_f$, the fixed point is
nonhyperbolic. Nevertheless, local identity still implies equality
of the corresponding center dynamics.

\begin{proposition}[Center-manifold invariance at the flip threshold]
\label{prop:center-manifold-invariance}
Assume \textbf{(H1)--(H3)} and set $\eta=\eta_f$.
For every fixed $c>0$, a sufficiently small local center manifold of
$G_{\eta_f}$ through $x_*$ is also a local center manifold of
$F_{\eta_f,c}$.

Moreover, the reduced maps on this common center manifold are
identical.
\end{proposition}

\begin{proof}
By Theorem~\ref{thm:local-identity}, the maps
$F_{\eta_f,c}$ and $G_{\eta_f}$ coincide in a neighborhood of
$x_*$. Choose a local center manifold of $G_{\eta_f}$ contained in
that neighborhood. Its local invariance equation is therefore also
satisfied by $F_{\eta_f,c}$. Since the restrictions of the two maps
to the manifold coincide pointwise, the corresponding reduced
one-dimensional dynamics are identical.
\end{proof}

\begin{remark}[Nonuniqueness of center manifolds]
\label{rem:center-manifold-nonuniqueness}
Center manifolds need not be unique as sets; the content of
Proposition~\ref{prop:center-manifold-invariance} is only that the
two maps admit one common choice with identical reduced dynamics.
\end{remark}

% ---------------------------------------------------------
\subsection{Consequences for the smooth-to-switching mechanism}
\label{subsec:local-structure-consequences}
% ---------------------------------------------------------

\begin{corollary}[Local bifurcation structure is clipping-independent]
\label{cor:local-bifurcation-clipping-independent}
Under \textbf{(H1)--(H4)}, for every fixed $c>0$:

\begin{enumerate}[label=\textnormal{(\roman*)}]

\item
the fixed point $x_*$ is unchanged by clipping;

\item
its first local stability threshold is
\[
\eta_f=\frac{2}{\lambda_*};
\]

\item
the critical eigendirection is $v_*$;

\item
the critical multiplier crosses $-1$ transversally at $\eta_f$;

\item
the local center-manifold reduction is identical to that of the
unclipped gradient map;

\item
the supercritical flip coefficient $\ell_*$ is independent of $c$.

\end{enumerate}
\end{corollary}

\begin{proof}
Items (i)--(vi) follow respectively from
Proposition~\ref{prop:fixed-point-invariance},
Theorem~\ref{thm:local-stability-threshold},
\eqref{eq:critical-eigenpair},
\eqref{eq:transversal-crossing},
Proposition~\ref{prop:center-manifold-invariance}, and
Corollary~\ref{cor:flip-coefficient-independence}.
\end{proof}

The key implication is that $c$ does not enter the local birth
mechanism of the period-doubled branch: its first dynamical effect
must occur at a distinct finite-amplitude event, when the
bifurcating orbit reaches the endogenous switching hypersurface
$\Sigma_c$.



% =========================================================
\section{Smooth flip bifurcation and the emerging period-two branch}
\label{sec:flip-bifurcation}
% =========================================================

In this section we study the smooth gradient map
\[
G_\eta(x)=x-\eta\nabla V(x)
\]
near the critical step size
\[
\eta_f=\frac{2}{\lambda_*}.
\]

The local mechanism considered here is the classical simple flip
bifurcation of a smooth discrete map
\cite{Kuznetsov2004,GuckenheimerHolmes1983}.
For gradient-descent dynamics, Mulayoff and Stich
\cite{MulayoffStich2026} recently derived a multivariate nonlinear
stability criterion involving the second-, third-, and fourth-order
derivatives of the potential, while Gan \cite{Gan2026} formulated
edge-of-stability dynamics in a systematic bifurcation-theoretic
framework.

Accordingly, the existence of a smooth period-two branch is not the
new contribution of the present paper.  We reproduce the local
reduction because its explicit coefficient enters the
switching-separation law derived later; we therefore give a
self-contained center-manifold and normal-form derivation in the
notation of the present paper.

By the local identity theorem established in Section~3, the
resulting local bifurcation structure is inherited exactly by the
clipped map $F_{\eta,c}$ as long as the bifurcating orbit remains
inside the smooth region $\Omega_c^-$.

Throughout this section we write
\begin{equation}
\mu
=
\eta-\eta_f.
\label{eq:flip-mu}
\end{equation}
Thus the bifurcation point corresponds to $\mu=0$.

Under \textbf{(H2)--(H3)}, the multiplier associated with $v_*$ is
\begin{equation}
\rho_*(\mu)
=
1-(\eta_f+\mu)\lambda_*
=
-1-\lambda_*\mu,
\label{eq:critical-multiplier-mu}
\end{equation}
whereas all remaining multipliers lie strictly inside the unit disk
at $\mu=0$.

The nonlinear character of the bifurcation is governed by the
coefficient
\[
\ell_*
=
\frac{\Gamma_*}{3\lambda_*},
\]
defined in \eqref{eq:ell-star}.  Assumption \textbf{(H4)} gives
\[
\ell_*>0,
\]
which corresponds to the supercritical case.

% ---------------------------------------------------------
\subsection{Parameter-dependent center reduction}
\label{subsec:parameter-center-reduction}
% ---------------------------------------------------------

Translate the critical point to the origin, $y=x-x_*$, and define
\begin{equation}
\Phi_\mu(y)
:=
G_{\eta_f+\mu}(x_*+y)-x_*.
\label{eq:Phi-mu}
\end{equation}
Then $\Phi_\mu(0)=0$ for every sufficiently small $\mu$, and at
$\mu=0$,
\begin{equation}
D\Phi_0(0)
=
I-\eta_fH_*,
\label{eq:Phi-linearization}
\end{equation}
whose unique critical multiplier is $-1$ with eigenvector $v_*$; the
complementary spectrum lies strictly inside the unit disk.

To incorporate the bifurcation parameter, consider the extended map
\begin{equation}
(y,\mu)
\longmapsto
\bigl(\Phi_\mu(y),\mu\bigr).
\label{eq:extended-map}
\end{equation}
Its two-dimensional local center manifold is associated with the
critical state direction $v_*$ and the neutral parameter direction,
and after restricting to a sufficiently small neighborhood can be
represented as
\begin{equation}
y
=
zv_*+h(z,\mu),
\label{eq:center-parametrization}
\end{equation}
with
\begin{equation}
h(z,\mu)\in v_*^\perp,
\qquad
h(0,\mu)=0,
\qquad
D_zh(0,0)=0,
\qquad
h(z,\mu)
=
O\bigl(z^2+|\mu||z|\bigr).
\label{eq:center-h-properties}
\end{equation}
The dynamics on the center manifold are described by the scalar map
\begin{equation}
z_{n+1}
=
f_\mu(z_n).
\label{eq:reduced-map}
\end{equation}

% ---------------------------------------------------------
\subsection{Flip normal form}
\label{subsec:flip-normal-form}
% ---------------------------------------------------------

We next record the local normal form required for the asymptotic
analysis.

\begin{proposition}[Reduced flip normal form]
\label{prop:flip-normal-form}
Assume \textbf{(H1)--(H4)}.
There exist local state and parameter coordinates, preserving
$\mu=0$ and $z=0$, in which the reduced map
\eqref{eq:reduced-map} takes the form
\begin{equation}
f_\mu(z)
=
-\bigl(1+\lambda_*\mu\bigr)z
+
\ell_*z^3
+
R(z,\mu),
\label{eq:flip-normal-form}
\end{equation}
where
\begin{equation}
R(z,\mu)
=
O\left(
|z|^4
+
|\mu||z|^2
+
\mu^2|z|
\right)
\label{eq:flip-remainder}
\end{equation}
as $(z,\mu)\to(0,0)$.

The cubic coefficient is
\begin{equation}
\ell_*
=
\frac{\eta_f}{6}
\left(
3T_*[v_*,v_*,H_*^{-1}g_*]
-
Q_*[v_*,v_*,v_*,v_*]
\right),
\label{eq:flip-coefficient-section3}
\end{equation}
or equivalently
\begin{equation}
\ell_*
=
\frac{\Gamma_*}{3\lambda_*}.
\label{eq:flip-coefficient-Gamma}
\end{equation}
\end{proposition}

\begin{proof}
The critical multiplier $-1$ is simple and crosses the unit circle
transversally (Section~\ref{subsec:spectral-consequences}), so the
standard parameter-dependent center-manifold and flip-normal-form
reductions apply; see \cite{Carr1981,HirschPughShub1977,Kuznetsov2004}.
The linear parameter coefficient is $\rho_*(\mu)=-1-\lambda_*\mu$
by \eqref{eq:critical-multiplier-mu}, and the quadratic term is
eliminated by the standard near-identity transformation for a simple
flip; the remaining cubic coefficient is the first flip coefficient.

For the gradient map
\[
G_\eta(x)=x-\eta\nabla V(x),
\]
one has
\begin{align}
D^2G_{\eta_f}(x_*)[u,w]
&=
-\eta_f
T_*[u,w,\cdot]^\sharp,
\label{eq:D2G}
\\
D^3G_{\eta_f}(x_*)[u,w,q]
&=
-\eta_f
Q_*[u,w,q,\cdot]^\sharp.
\label{eq:D3G}
\end{align}
The standard flip-coefficient formula therefore gives
\[
\ell_*
=
\frac{\eta_f}{2}
T_*[v_*,v_*,H_*^{-1}g_*]
-
\frac{\eta_f}{6}
Q_*[v_*,v_*,v_*,v_*].
\]
Using
\[
\eta_f=\frac{2}{\lambda_*}
\]
and the definition of $\Gamma_*$ in
\eqref{eq:Gamma-star} yields
\[
\ell_*
=
\frac{\Gamma_*}{3\lambda_*}.
\]

The remainder estimate
\eqref{eq:flip-remainder} follows from the regularity in
\textbf{(H1)} and the parameter-dependent near-identity normal-form
transformation.
\end{proof}

\begin{remark}
\label{rem:flip-sign-convention}
The sign convention in \eqref{eq:flip-normal-form} is fixed
throughout the paper:
\[
\ell_*>0
\]
corresponds to a supercritical flip bifurcation, with a stable
period-two branch emerging for $\mu>0$.
This convention is consistent with the multivariate
gradient-descent flip criterion encoded by \textbf{(H4)}.
\end{remark}

% ---------------------------------------------------------
\subsection{The second iterate}
\label{subsec:second-iterate}
% ---------------------------------------------------------

The bifurcating period-two orbit is most conveniently obtained from
the fixed points of the second iterate.

\begin{lemma}[Expansion of the second iterate]
\label{lem:second-iterate-expansion}
Under the hypotheses of
Proposition~\ref{prop:flip-normal-form},
\begin{equation}
f_\mu^2(z)-z
=
2\lambda_*\mu z
-
2\ell_*z^3
+
\mathcal R_2(z,\mu),
\label{eq:second-iterate-expansion}
\end{equation}
where
\begin{equation}
\mathcal R_2(z,\mu)
=
O\left(
|z|^4
+
|\mu||z|^2
+
\mu^2|z|
\right).
\label{eq:second-iterate-remainder}
\end{equation}
\end{lemma}

\begin{proof}
Write
\[
a_\mu:=1+\lambda_*\mu.
\]
Then
\[
f_\mu(z)
=
-a_\mu z+\ell_*z^3+R(z,\mu).
\]
Ignoring terms of order contained in
\eqref{eq:second-iterate-remainder}, one obtains
\begin{align}
f_\mu^2(z)
&=
-a_\mu
\left(
-a_\mu z+\ell_*z^3
\right)
+
\ell_*
\left(
-a_\mu z
\right)^3
+
\mathcal R_2(z,\mu)
\nonumber\\
&=
a_\mu^2z
-
\ell_*a_\mu(1+a_\mu^2)z^3
+
\mathcal R_2(z,\mu).
\label{eq:second-compose}
\end{align}
Since
\[
a_\mu^2
=
1+2\lambda_*\mu+O(\mu^2)
\]
and
\[
a_\mu(1+a_\mu^2)
=
2+O(\mu),
\]
we obtain
\[
f_\mu^2(z)-z
=
2\lambda_*\mu z
-
2\ell_*z^3
+
\mathcal R_2(z,\mu).
\]
The stated remainder estimate follows by substituting
\eqref{eq:flip-remainder} into the composition.
\end{proof}

Equation \eqref{eq:second-iterate-expansion} exhibits explicitly the
balance responsible for the square-root amplitude law:
\[
\lambda_*\mu
\sim
\ell_*z^2.
\]

% ---------------------------------------------------------
\subsection{Existence and amplitude of the period-two branch}
\label{subsec:period-two-amplitude}
% ---------------------------------------------------------

We now obtain the local period-two branch together with its leading
amplitude.

\begin{theorem}[Supercritical period-two branch]
\label{thm:period-two-branch}
Assume \textbf{(H1)--(H4)}.
There exists $\mu_0>0$ such that for every
\[
0<\mu<\mu_0,
\]
the smooth gradient map $G_{\eta_f+\mu}$ possesses a unique local
period-two cycle
\begin{equation}
\mathcal O_2(\mu)
=
\left\{
x_-(\mu),x_+(\mu)
\right\}
\label{eq:period-two-cycle}
\end{equation}
in a sufficiently small neighborhood of $x_*$.

The two points satisfy
\begin{equation}
x_\pm(\mu)
=
x_*
\pm
A_*\mu^{1/2}v_*
+
O(\mu),
\qquad
\mu\to0^+,
\label{eq:period-two-x-expansion}
\end{equation}
where
\begin{equation}
A_*
=
\sqrt{\frac{\lambda_*}{\ell_*}}
\label{eq:A-star-ell}
\end{equation}
and therefore
\begin{equation}
A_*^2
=
\frac{3\lambda_*^2}{\Gamma_*}.
\label{eq:A-star-Gamma}
\end{equation}

The labeling may be chosen so that
\begin{equation}
G_{\eta_f+\mu}(x_+(\mu))
=
x_-(\mu),
\qquad
G_{\eta_f+\mu}(x_-(\mu))
=
x_+(\mu).
\label{eq:period-two-exchange}
\end{equation}
\end{theorem}

\begin{proof}
We first work in the reduced coordinate $z$.

Set
\[
\mu=\varepsilon^2,
\qquad
\varepsilon>0,
\]
and seek nonzero solutions of
\[
f_{\varepsilon^2}^2(z)-z=0
\]
in the form
\[
z=\varepsilon w.
\]
By Lemma~\ref{lem:second-iterate-expansion},
\begin{align}
f_{\varepsilon^2}^2(\varepsilon w)-\varepsilon w
&=
2\lambda_*\varepsilon^3w
-
2\ell_*\varepsilon^3w^3
+
O(\varepsilon^4)
\nonumber\\
&=
2\varepsilon^3
\left[
w(\lambda_*-\ell_*w^2)
+
O(\varepsilon)
\right].
\label{eq:scaled-second-iterate}
\end{align}

Define the rescaled bifurcation equation
\begin{equation}
\mathscr F(w,\varepsilon)
:=
\frac{
f_{\varepsilon^2}^2(\varepsilon w)-\varepsilon w
}{
2\varepsilon^3
},
\qquad
\varepsilon>0.
\label{eq:rescaled-F}
\end{equation}
Under \textbf{(H1)} the rescaled map extends as a $C^3$ function to
$\varepsilon=0$: since $V\in C^7$, the gradient map $G_\eta$ is
$C^6$, and by the parameter-dependent flip reduction so is the
reduced map $f_\mu$; the numerator in \eqref{eq:rescaled-F} is
therefore $C^6$ in $(\varepsilon,w)$ and vanishes to third order in
$\varepsilon$, so dividing by $2\varepsilon^3$ yields a $C^3$
function.  We set
\begin{equation}
\mathscr F(w,0)
=
w(\lambda_*-\ell_*w^2).
\label{eq:rescaled-limit}
\end{equation}

Since $\ell_*>0$, the nonzero zeros of
$\mathscr F(\cdot,0)$ are
\begin{equation}
w_\pm^0
=
\pm
\sqrt{\frac{\lambda_*}{\ell_*}}.
\label{eq:w0}
\end{equation}
Furthermore,
\begin{align}
\partial_w\mathscr F(w_\pm^0,0)
&=
\lambda_*
-
3\ell_*(w_\pm^0)^2
\nonumber\\
&=
-2\lambda_*
\neq0.
\label{eq:IFT-nondegenerate}
\end{align}
The implicit function theorem therefore yields two branches
\[
w_\pm(\varepsilon)
=
w_\pm^0+O(\varepsilon).
\]
Consequently,
\begin{align}
z_\pm(\mu)
&=
\varepsilon w_\pm(\varepsilon)
\nonumber\\
&=
\pm
\sqrt{\frac{\lambda_*}{\ell_*}}
\sqrt{\mu}
+
O(\mu).
\label{eq:z-period-two}
\end{align}

The fixed point $z=0$ is locally the only fixed point of
$f_\mu$ near the origin.  Indeed,
\[
\partial_z(f_0(z)-z)\big|_{z=0}
=
-2
\neq0,
\]
so the implicit function theorem excludes nearby nonzero
period-one fixed points.  Hence the two nonzero solutions
$z_\pm(\mu)$ of
\[
f_\mu^2(z)=z
\]
have prime period two.

Returning to the original coordinates through
\eqref{eq:center-parametrization},
\[
x_\pm(\mu)
=
x_*
+
z_\pm(\mu)v_*
+
h(z_\pm(\mu),\mu).
\]
Since
\[
h(z,\mu)
=
O(z^2+|\mu||z|)
\]
and
\[
z_\pm(\mu)=O(\sqrt{\mu}),
\]
we have
\[
h(z_\pm(\mu),\mu)=O(\mu).
\]
Therefore
\[
x_\pm(\mu)
=
x_*
\pm
\sqrt{\frac{\lambda_*}{\ell_*}}
\sqrt{\mu}\,v_*
+
O(\mu),
\]
which proves \eqref{eq:period-two-x-expansion}.

Finally, the two nonzero points form a single two-cycle, so after
possibly interchanging the labels we obtain
\eqref{eq:period-two-exchange}.

The uniqueness of the sufficiently small period-two cycle follows
from the local center-manifold reduction and the uniqueness of the
two nonzero solution branches of the reduced second-iterate
equation.
\end{proof}

% ---------------------------------------------------------
\subsection{Stability of the emerging cycle}
\label{subsec:period-two-stability}
% ---------------------------------------------------------

We next verify the local stability of the branch.

\begin{theorem}[Stability of the supercritical period-two cycle]
\label{thm:period-two-stability}
Under \textbf{(H1)--(H4)}, the period-two cycle
$\mathcal O_2(\mu)$ given by
Theorem~\ref{thm:period-two-branch} is locally asymptotically stable
for all sufficiently small $\mu>0$.

Along the critical center direction, the multiplier of the
two-step map satisfies
\begin{equation}
\rho_{\mathrm{cyc}}^{c}(\mu)
=
1-4\lambda_*\mu
+
O(\mu^{3/2}),
\qquad
\mu\to0^+.
\label{eq:critical-cycle-multiplier}
\end{equation}
\end{theorem}

\begin{proof}
Differentiate the expansion
\eqref{eq:second-iterate-expansion}:
\begin{equation}
(f_\mu^2)'(z)
=
1
+
2\lambda_*\mu
-
6\ell_*z^2
+
O\left(
|z|^3
+
|\mu||z|
+
\mu^2
\right).
\label{eq:second-iterate-derivative}
\end{equation}
At either point $z=z_\pm(\mu)$,
Theorem~\ref{thm:period-two-branch} gives
\[
z_\pm(\mu)^2
=
\frac{\lambda_*}{\ell_*}\mu
+
O(\mu^{3/2}).
\]
Hence
\begin{align}
(f_\mu^2)'(z_\pm(\mu))
&=
1
+
2\lambda_*\mu
-
6\lambda_*\mu
+
O(\mu^{3/2})
\nonumber\\
&=
1
-
4\lambda_*\mu
+
O(\mu^{3/2}).
\end{align}
Thus the center multiplier lies strictly inside the unit disk for
all sufficiently small $\mu>0$.

At $\mu=0$, every transverse multiplier of the two-step map equals
\[
\rho_i(\eta_f)^2,
\qquad
i=1,\ldots,d-1,
\]
and by \eqref{eq:noncritical-spectrum},
\[
|\rho_i(\eta_f)|<1.
\]
Therefore
\[
\rho_i(\eta_f)^2<1.
\]
By continuity, all transverse multipliers remain strictly inside the
unit disk along the sufficiently small period-two branch.

The reduction principle then implies local asymptotic stability of
the full period-two cycle.
\end{proof}

\begin{remark}
\label{rem:cycle-multiplier-meaning}
Equation \eqref{eq:critical-cycle-multiplier} shows that the
newly created cycle is only weakly attracting immediately after the
flip:
\[
\rho_{\mathrm{cyc}}^{c}(\mu)\to1^-
\qquad
\text{as }\mu\to0^+.
\]
This weak attraction is relevant later when the branch approaches
the switching hypersurface, because the switching interaction occurs
on a parameter scale that also tends to zero with $c$.
\end{remark}

% ---------------------------------------------------------
\subsection{Transfer to the clipped map before switching}
\label{subsec:transfer-clipped}
% ---------------------------------------------------------

We finally transfer the preceding smooth bifurcation result to the
clipped map.

\begin{corollary}[Pre-contact period-two cycle for the clipped map]
\label{cor:precontact-clipped-cycle}
Assume \textbf{(H1)--(H5)} and fix $c>0$.
Then there exists
\[
\mu_c>0
\]
such that for every
\[
0<\mu<\mu_c,
\]
the clipped map
\[
F_{\eta_f+\mu,c}
\]
possesses the same locally asymptotically stable period-two cycle
$\mathcal O_2(\mu)$ as the smooth gradient map
$G_{\eta_f+\mu}$.

Moreover,
\begin{equation}
x_\pm(\mu)
\in
\Omega_c^-,
\qquad
0<\mu<\mu_c.
\label{eq:cycle-inside-smooth-region}
\end{equation}
\end{corollary}

\begin{proof}
For the fixed value $c>0$, let $r_c>0$ be given by
Theorem~\ref{thm:local-identity}. By
Theorem~\ref{thm:period-two-branch},
\[
x_\pm(\mu)\to x_*
\qquad
\text{as }\mu\to0^+.
\]
Hence there exists $\mu_c>0$ such that
\[
x_\pm(\mu)\in B_{r_c}(x_*)
\subset\Omega_c^-
\]
for every $0<\mu<\mu_c$.

On this ball,
\[
F_{\eta_f+\mu,c}
=
G_{\eta_f+\mu}
\]
by Theorem~\ref{thm:local-identity}. Therefore the two maps have
exactly the same period-two cycle and the same local derivatives
along that cycle.  Its local asymptotic stability follows from
Theorem~\ref{thm:period-two-stability}.
\end{proof}

\begin{remark}[The first possible effect of clipping]
\label{rem:first-clipping-effect}
Corollary~\ref{cor:precontact-clipped-cycle} shows that the
period-two branch is born entirely inside the smooth region.
Therefore clipping does not influence either its creation or its
initial stability.

The first possible effect of clipping occurs only when one of the
points $x_\pm(\mu)$ reaches the endogenous switching hypersurface
\[
\Sigma_c
=
\left\{
x:
\|\nabla V(x)\|_\diamond=c
\right\}.
\]
Determining the parameter value of this first contact is the central
problem of Sections~\ref{sec:switching-geometry} and
\ref{sec:separation-law}.
\end{remark}



% =========================================================
\section{Switching geometry along the period-two branch}
\label{sec:switching-geometry}
% =========================================================

We now study how the smooth period-two branch constructed in
Section~\ref{sec:flip-bifurcation} approaches the endogenous
switching hypersurface
\[
\Sigma_c
=
\left\{
x\in\mathcal U:
\|\nabla V(x)\|_\diamond=c
\right\}.
\]

A key structural feature of gradient maps appears at this point.
For an arbitrary period-two cycle of gradient descent, the gradients
at its two points are exactly opposite. Consequently, their gradient
norms are identical for every norm $\|\cdot\|_\diamond$.

This identity has two important consequences.

First, the two points of the smooth period-two orbit necessarily
reach the clipping boundary simultaneously. Thus, unlike a generic
piecewise-smooth map, the present gradient-generated switching
problem does not exhibit a splitting of the two first-contact
parameters.

Second, the common gradient magnitude along the bifurcating cycle
has a particularly simple expansion in the square-root parameter
$\varepsilon=\sqrt{\eta-\eta_f}$. This structure will yield the
quadratic separation law in Section~\ref{sec:separation-law}.

% ---------------------------------------------------------
\subsection{Exact gradient balance on a period-two cycle}
\label{subsec:exact-gradient-balance}
% ---------------------------------------------------------

We begin with an exact identity that does not rely on the local
bifurcation expansion. An equivalent balance relation appears in
the recent characterization of gradient-descent period-two orbits
by Litman \cite{Litman2026}. Since this identity is the starting
point of the switching analysis below, we include its elementary
proof for completeness.

\begin{proposition}[Exact gradient balance]
\label{prop:exact-gradient-balance}
Let $\eta>0$, and suppose that
$\{x_-,x_+\}$ is a period-two cycle of the smooth gradient map
\[
G_\eta(x)
=
x-\eta\nabla V(x),
\]
that is,
\begin{equation}
G_\eta(x_+)=x_-,
\qquad
G_\eta(x_-)=x_+,
\qquad
x_+\neq x_-.
\label{eq:two-cycle-relations}
\end{equation}
Then
\begin{equation}
\nabla V(x_+)
=
\frac{x_+-x_-}{\eta},
\qquad
\nabla V(x_-)
=
-\frac{x_+-x_-}{\eta}.
\label{eq:exact-gradient-opposition}
\end{equation}
In particular,
\begin{equation}
\nabla V(x_+)
=
-\nabla V(x_-).
\label{eq:opposite-gradients}
\end{equation}
Consequently, for every norm $\|\cdot\|_\diamond$,
\begin{equation}
\|\nabla V(x_+)\|_\diamond
=
\|\nabla V(x_-)\|_\diamond.
\label{eq:equal-gradient-norms}
\end{equation}
\end{proposition}

\begin{proof}
The first identity in \eqref{eq:two-cycle-relations} gives
\[
x_-
=
x_+
-
\eta\nabla V(x_+),
\]
and hence
\[
\nabla V(x_+)
=
\frac{x_+-x_-}{\eta}.
\]
Similarly,
\[
x_+
=
x_-
-
\eta\nabla V(x_-)
\]
implies
\[
\nabla V(x_-)
=
\frac{x_--x_+}{\eta}
=
-\frac{x_+-x_-}{\eta}.
\]
This proves \eqref{eq:exact-gradient-opposition} and
\eqref{eq:opposite-gradients}.

Since every norm is even,
\[
\|-g\|_\diamond=\|g\|_\diamond,
\]
and therefore
\[
\|\nabla V(x_-)\|_\diamond
=
\|-\nabla V(x_+)\|_\diamond
=
\|\nabla V(x_+)\|_\diamond.
\]
\end{proof}

\begin{remark}[Gradient structure versus a generic switching map]
\label{rem:gradient-special-structure}
Equation \eqref{eq:equal-gradient-norms} is a special consequence of
the gradient-descent structure.  In a generic piecewise-smooth map,
the two points of a period-two orbit need not have equal switching
functions and may therefore encounter a switching boundary at
different parameter values.

For radial gradient clipping, by contrast, the switching function
depends only on
\[
\|\nabla V(x)\|_\diamond,
\]
and Proposition~\ref{prop:exact-gradient-balance} forces the two
switching values to be identical along every smooth period-two
cycle.
\end{remark}

% ---------------------------------------------------------
\subsection{Midpoint and half-difference coordinates}
\label{subsec:midpoint-difference}
% ---------------------------------------------------------

For the bifurcating period-two cycle
\[
\mathcal O_2(\mu)
=
\{x_-(\mu),x_+(\mu)\},
\qquad
\mu=\eta-\eta_f>0,
\]
define its midpoint and half-difference by
\begin{align}
m(\mu)
&:=
\frac{x_+(\mu)+x_-(\mu)}{2}-x_*,
\label{eq:midpoint-definition}
\\
d(\mu)
&:=
\frac{x_+(\mu)-x_-(\mu)}{2}.
\label{eq:half-difference-definition}
\end{align}
Thus
\begin{equation}
x_\pm(\mu)
=
x_*+m(\mu)\pm d(\mu).
\label{eq:midpoint-difference-representation}
\end{equation}

Theorem~\ref{thm:period-two-branch} implies
\begin{equation}
m(\mu)=O(\mu),
\qquad
d(\mu)
=
A_*\sqrt{\mu}\,v_*
+
O(\mu).
\label{eq:rough-md-orders}
\end{equation}

The period-two equations yield considerably more structure.

\begin{lemma}[Midpoint--difference equations]
\label{lem:midpoint-difference-equations}
For every smooth period-two cycle of $G_\eta$ written in the form
\eqref{eq:midpoint-difference-representation},
\begin{align}
&
\nabla V(x_*+m+d)
+
\nabla V(x_*+m-d)
=
0,
\label{eq:midpoint-equation}
\\
&
\nabla V(x_*+m+d)
-
\nabla V(x_*+m-d)
=
\frac{4}{\eta}d.
\label{eq:difference-equation}
\end{align}
\end{lemma}

\begin{proof}
By Proposition~\ref{prop:exact-gradient-balance},
\[
\nabla V(x_*+m+d)
=
\frac{2}{\eta}d
\]
and
\[
\nabla V(x_*+m-d)
=
-\frac{2}{\eta}d.
\]
Adding and subtracting these identities gives
\eqref{eq:midpoint-equation} and
\eqref{eq:difference-equation}.
\end{proof}

% ---------------------------------------------------------
\subsection{Square-root parametrization and parity}
\label{subsec:square-root-parity}
% ---------------------------------------------------------

By the parameter-dependent flip theorem
\cite{Kuznetsov2004,GuckenheimerHolmes1983} and \textbf{(H1)}, the
local period-two branch admits a $C^4$ signed parametrization in the
variable $\varepsilon$ introduced below; in particular the midpoint
and half-difference functions are $C^4$, and the parity relations
derived next are justified to the orders used in
Proposition~\ref{prop:refined-cycle-expansion}.

Introduce the signed square-root parameter
\begin{equation}
\varepsilon^2=\mu,
\qquad
\eta=\eta_f+\varepsilon^2.
\label{eq:epsilon-definition}
\end{equation}

The two-cycle equations are invariant under the transformation
\[
(d,\varepsilon)
\longmapsto
(-d,-\varepsilon),
\]
which merely exchanges the two points of the cycle.

The local uniqueness of the bifurcating period-two branch therefore
allows the branch to be parametrized so that
\begin{equation}
m(-\varepsilon)=m(\varepsilon),
\qquad
d(-\varepsilon)=-d(\varepsilon).
\label{eq:parity-md}
\end{equation}

Hence the midpoint is even and the half-difference is odd in
$\varepsilon$.

\begin{proposition}[Refined two-cycle expansion]
\label{prop:refined-cycle-expansion}
Assume \textbf{(H1)--(H4)}.
The period-two branch can be parametrized by the signed variable
$\varepsilon$, with $\mu=\varepsilon^2$, such that
\begin{align}
m(\varepsilon)
&=
\varepsilon^2m_2
+
O(\varepsilon^4),
\label{eq:m-expansion}
\\
d(\varepsilon)
&=
A_*\varepsilon v_*
+
O(\varepsilon^3),
\label{eq:d-expansion}
\end{align}
where
\begin{equation}
m_2
=
-\frac{A_*^2}{2}H_*^{-1}g_*
=
-\frac{A_*^2}{2}h_*.
\label{eq:m2-formula}
\end{equation}

Consequently,
\begin{equation}
x_\pm(\mu)
=
x_*
-
\frac{A_*^2}{2}\mu h_*
\pm
A_*\sqrt{\mu}\,v_*
+
O(\mu^{3/2}).
\label{eq:refined-xpm-expansion}
\end{equation}
\end{proposition}

\begin{proof}
The parity relations \eqref{eq:parity-md} imply that the local
expansions have the form
\[
m(\varepsilon)
=
\varepsilon^2m_2+O(\varepsilon^4)
\]
and
\[
d(\varepsilon)
=
\varepsilon d_1+O(\varepsilon^3).
\]
By Theorem~\ref{thm:period-two-branch},
\[
d_1=A_*v_*.
\]

It remains to determine $m_2$.
Taylor expansion of the gradient about $x_*$ gives
\begin{equation}
\nabla V(x_*+y)
=
H_*y
+
\frac12
T_*[y,y,\cdot]^\sharp
+
O(\|y\|_2^3).
\label{eq:gradient-second-order-Taylor}
\end{equation}

Apply \eqref{eq:gradient-second-order-Taylor} to
$y=m+d$ and $y=m-d$ and add the resulting expressions.
Using \eqref{eq:midpoint-equation}, we obtain
\begin{equation}
2H_*m
+
T_*[d,d,\cdot]^\sharp
+
O(\varepsilon^4)
=
0.
\label{eq:midpoint-expanded}
\end{equation}
Here we used
\[
m=O(\varepsilon^2),
\qquad
d=O(\varepsilon),
\]
and the cancellation of all terms odd in $d$.

Substituting
\[
m=\varepsilon^2m_2+O(\varepsilon^4)
\]
and
\[
d=A_*\varepsilon v_*+O(\varepsilon^3)
\]
into \eqref{eq:midpoint-expanded} gives
\[
2H_*m_2
+
A_*^2
T_*[v_*,v_*,\cdot]^\sharp
=
0.
\]
Since $H_*$ is invertible,
\[
m_2
=
-\frac{A_*^2}{2}
H_*^{-1}
T_*[v_*,v_*,\cdot]^\sharp
=
-\frac{A_*^2}{2}h_*.
\]
This proves \eqref{eq:m2-formula} and hence
\eqref{eq:refined-xpm-expansion}.
\end{proof}

\begin{remark}[Nonlinear displacement of the cycle midpoint]
\label{rem:midpoint-drift}
Although the leading oscillation is aligned with the maximal
Hessian eigendirection $v_*$, the midpoint of the bifurcating
two-cycle generally does not remain at $x_*$. Instead,
\[
\frac{x_+(\mu)+x_-(\mu)}{2}
=
x_*
-
\frac{A_*^2}{2}\mu h_*
+
O(\mu^2).
\]
Thus the third derivative of $V$ produces an $O(\mu)$ drift of the
cycle center through the vector
\[
h_*
=
H_*^{-1}
T_*[v_*,v_*,\cdot]^\sharp.
\]
\end{remark}

% ---------------------------------------------------------
\subsection{A consistency derivation of the flip amplitude}
\label{subsec:amplitude-consistency}
% ---------------------------------------------------------

The midpoint--difference equations also provide a direct,
coordinate-free derivation of the leading oscillation amplitude.

\begin{proposition}[Amplitude identity]
\label{prop:amplitude-identity}
Under \textbf{(H1)--(H4)}, the leading amplitude satisfies
\begin{equation}
A_*^2\Gamma_*
=
3\lambda_*^2.
\label{eq:amplitude-Gamma-identity}
\end{equation}
Equivalently,
\begin{equation}
A_*^2
=
\frac{3\lambda_*^2}{\Gamma_*},
\label{eq:amplitude-coordinate-free}
\end{equation}
in agreement with
\eqref{eq:A-star-Gamma}.
\end{proposition}

\begin{proof}
Expand the difference equation
\eqref{eq:difference-equation} to cubic order.

Using
\[
m
=
\varepsilon^2m_2+O(\varepsilon^4),
\qquad
d
=
A_*\varepsilon v_*+O(\varepsilon^3),
\]
Taylor expansion gives
\begin{align}
&
\nabla V(x_*+m+d)
-
\nabla V(x_*+m-d)
\nonumber\\
&\quad=
2H_*d
+
2T_*[m,d,\cdot]^\sharp
+
\frac13Q_*[d,d,d,\cdot]^\sharp
+
O(\varepsilon^4).
\label{eq:difference-expanded}
\end{align}

On the other hand,
\[
\frac{4}{\eta}
=
\frac{4}{\eta_f+\varepsilon^2}
=
\frac{4}{\eta_f}
-
\frac{4}{\eta_f^2}\varepsilon^2
+
O(\varepsilon^4).
\label{eq:four-over-eta}
\]

Take the Euclidean inner product of
\eqref{eq:difference-equation} with $v_*$ and compare the
coefficients of $\varepsilon^3$.

The unknown $O(\varepsilon^3)$ correction in $d$ cancels because
\[
\frac{4}{\eta_f}
=
2\lambda_*.
\]
Thus,
\begin{equation}
2A_*
T_*[v_*,v_*,m_2]
+
\frac{A_*^3}{3}
Q_*[v_*,v_*,v_*,v_*]
=
-\frac{4A_*}{\eta_f^2}.
\label{eq:amplitude-before-substitution}
\end{equation}

Since
\[
m_2
=
-\frac{A_*^2}{2}h_*
\]
and
\[
\eta_f=\frac{2}{\lambda_*},
\]
equation \eqref{eq:amplitude-before-substitution} becomes
\[
-A_*^3T_*[v_*,v_*,h_*]
+
\frac{A_*^3}{3}Q_*[v_*^4]
=
-A_*\lambda_*^2.
\]
Dividing by $A_*>0$ and rearranging gives
\[
A_*^2
\left(
3T_*[v_*,v_*,h_*]
-
Q_*[v_*^4]
\right)
=
3\lambda_*^2.
\]
By definition of $\Gamma_*$,
\[
A_*^2\Gamma_*
=
3\lambda_*^2.
\]
\end{proof}

\begin{remark}
\label{rem:coordinate-free-check}
Proposition~\ref{prop:amplitude-identity} provides an independent
check of the center-manifold normal-form calculation in
Section~\ref{sec:flip-bifurcation}.  In particular, the same
high-order quantity $\Gamma_*$ governs both the supercriticality
criterion and the leading amplitude of the emerging period-two
cycle.
\end{remark}

% ---------------------------------------------------------
\subsection{Gradient exposure along the bifurcating cycle}
\label{subsec:gradient-exposure}
% ---------------------------------------------------------

Define the common gradient exposure of the smooth period-two cycle
by
\begin{equation}
\mathscr E(\mu)
:=
\|\nabla V(x_+(\mu))\|_\diamond
=
\|\nabla V(x_-(\mu))\|_\diamond.
\label{eq:exposure-definition}
\end{equation}
The equality is exact by
Proposition~\ref{prop:exact-gradient-balance}.

Using the half-difference coordinate, one obtains the exact formula
\begin{equation}
\mathscr E(\mu)
=
\frac{2}{\eta_f+\mu}
\|d(\mu)\|_\diamond.
\label{eq:exact-exposure-d}
\end{equation}

This identity is considerably stronger than a direct Taylor
expansion of $\nabla V$ at the two cycle points.

\begin{theorem}[Asymptotic switching geometry of the two-cycle]
\label{thm:exposure-expansion}
Assume \textbf{(H1)--(H5)}.
As $\mu\to0^+$,
\begin{equation}
\mathscr E(\mu)
=
B_*\mu^{1/2}
+
O(\mu^{3/2}),
\label{eq:exposure-expansion}
\end{equation}
where
\begin{equation}
B_*
=
\lambda_*A_*
\|v_*\|_\diamond.
\label{eq:B-star}
\end{equation}

Equivalently,
\begin{equation}
B_*^2
=
\frac{
3\lambda_*^4
\|v_*\|_\diamond^2
}{
\Gamma_*
}.
\label{eq:B-star-Gamma}
\end{equation}

In particular, there is no term of order $\mu$ in
\eqref{eq:exposure-expansion}.
\end{theorem}

\begin{proof}
Set
\[
\varepsilon=\sqrt{\mu}.
\]
By Proposition~\ref{prop:refined-cycle-expansion},
\[
d(\varepsilon)
=
A_*\varepsilon v_*
+
O(\varepsilon^3).
\]
Using positive homogeneity of the norm,
\begin{align}
\|d(\varepsilon)\|_\diamond
&=
\varepsilon
\left\|
A_*v_*
+
O(\varepsilon^2)
\right\|_\diamond
\nonumber\\
&=
A_*
\|v_*\|_\diamond
\varepsilon
+
O(\varepsilon^3).
\label{eq:d-norm-expansion}
\end{align}
Furthermore,
\begin{equation}
\frac{2}{\eta_f+\varepsilon^2}
=
\frac{2}{\eta_f}
+
O(\varepsilon^2)
=
\lambda_*
+
O(\varepsilon^2).
\label{eq:two-over-eta-expansion}
\end{equation}
Substituting
\eqref{eq:d-norm-expansion} and
\eqref{eq:two-over-eta-expansion} into
\eqref{eq:exact-exposure-d} yields
\[
\mathscr E(\mu)
=
\lambda_*A_*
\|v_*\|_\diamond
\varepsilon
+
O(\varepsilon^3).
\]
Since $\varepsilon=\sqrt{\mu}$,
\[
\mathscr E(\mu)
=
B_*\mu^{1/2}
+
O(\mu^{3/2}),
\]
where
\[
B_*
=
\lambda_*A_*
\|v_*\|_\diamond.
\]

Finally, using
\[
A_*^2
=
\frac{3\lambda_*^2}{\Gamma_*}
\]
gives
\[
B_*^2
=
\lambda_*^2
A_*^2
\|v_*\|_\diamond^2
=
\frac{
3\lambda_*^4
\|v_*\|_\diamond^2
}{
\Gamma_*
}.
\]
\end{proof}

\begin{remark}[Absence of a quadratic term in the exposure]
\label{rem:no-mu-term}
A naive substitution of
\[
x_\pm(\mu)
=
x_*
\pm
A_*\sqrt{\mu}\,v_*
+
O(\mu)
\]
into the gradient norm might suggest an expansion of the form
\[
B_1\sqrt{\mu}+B_2\mu+\cdots.
\]
The exact two-cycle identity
\eqref{eq:exact-exposure-d} shows that this is not the case.

The half-difference $d$ is odd in the signed square-root parameter
$\varepsilon$, while $\eta=\eta_f+\varepsilon^2$ is even.
Consequently, the common exposure is an odd function of
$\varepsilon$ to leading orders:
\[
\mathscr E(\varepsilon^2)
=
B_*\varepsilon
+
O(\varepsilon^3).
\]
This structural cancellation will improve the remainder in the
switching-contact separation law from the naively expected
$O(c^3)$ to $O(c^4)$.
\end{remark}

% ---------------------------------------------------------
\subsection{Simultaneous first contact}
\label{subsec:simultaneous-contact}
% ---------------------------------------------------------

We now record the most important geometric consequence of the exact
gradient balance.

\begin{theorem}[Simultaneous switching contact]
\label{thm:simultaneous-contact}
Let $\mathcal O_2(\mu)=\{x_-(\mu),x_+(\mu)\}$ be the smooth
period-two branch from
Theorem~\ref{thm:period-two-branch}.
For every sufficiently small $\mu>0$,
\begin{equation}
x_+(\mu)\in\Sigma_c
\quad\Longleftrightarrow\quad
x_-(\mu)\in\Sigma_c.
\label{eq:simultaneous-contact-equivalence}
\end{equation}

Likewise,
\begin{align}
x_+(\mu)\in\Omega_c^-
&\quad\Longleftrightarrow\quad
x_-(\mu)\in\Omega_c^-,
\label{eq:simultaneous-smooth}
\\
x_+(\mu)\in\Omega_c^+
&\quad\Longleftrightarrow\quad
x_-(\mu)\in\Omega_c^+.
\label{eq:simultaneous-clipped}
\end{align}

Hence the two points of a smooth gradient-descent period-two cycle
cannot cross a radial clipping boundary sequentially.
\end{theorem}

\begin{proof}
By Proposition~\ref{prop:exact-gradient-balance},
\[
\|\nabla V(x_+(\mu))\|_\diamond
=
\|\nabla V(x_-(\mu))\|_\diamond.
\]
Each of the three statements follows immediately by comparing this
common value with the clipping threshold $c$.
\end{proof}

\begin{corollary}[No pre-contact splitting of the two-cycle]
\label{cor:no-contact-splitting}
There do not exist two distinct first-contact parameters
\[
\eta_{\mathrm{sc}}^+(c)
\neq
\eta_{\mathrm{sc}}^-(c)
\]
for the two points of the smooth period-two branch.
Whenever the first switching contact exists, it is a simultaneous
two-point contact.
\end{corollary}

\begin{proof}
If one point lies on $\Sigma_c$, then the other lies on $\Sigma_c$
by Theorem~\ref{thm:simultaneous-contact}.
\end{proof}

% ---------------------------------------------------------
\subsection{Impossibility of a mixed smooth--clipped two-cycle}
\label{subsec:no-mixed-two-cycle}
% ---------------------------------------------------------

The simultaneous-contact structure is not restricted to the smooth
branch before contact.  Radial clipping itself rules out a genuine
period-two cycle having one point in the smooth region and the other
in the clipped region.

\begin{proposition}[No mixed-itinerary period-two cycle]
\label{prop:no-mixed-period-two}
Let
\[
\{y_-,y_+\}
\]
be a period-two cycle of the clipped map
\[
F_{\eta,c}(x)
=
x-\eta
\mathcal C_c^\diamond(\nabla V(x)).
\]
Then it is impossible to have
\begin{equation}
y_-\in\Omega_c^-,
\qquad
y_+\in\Omega_c^+,
\label{eq:forbidden-mixed-cycle}
\end{equation}
or vice versa.

Therefore every period-two cycle of $F_{\eta,c}$ is either entirely
smooth, entirely clipped, or lies on the switching boundary.
\end{proposition}

\begin{proof}
The period-two relations give
\begin{align}
\mathcal C_c^\diamond(\nabla V(y_+))
&=
\frac{y_+-y_-}{\eta},
\label{eq:clipped-balance-plus}
\\
\mathcal C_c^\diamond(\nabla V(y_-))
&=
-\frac{y_+-y_-}{\eta}.
\label{eq:clipped-balance-minus}
\end{align}
Hence
\begin{equation}
\left\|
\mathcal C_c^\diamond(\nabla V(y_+))
\right\|_\diamond
=
\left\|
\mathcal C_c^\diamond(\nabla V(y_-))
\right\|_\diamond.
\label{eq:clipped-output-equal}
\end{equation}

For radial clipping,
\begin{equation}
\left\|
\mathcal C_c^\diamond(g)
\right\|_\diamond
=
\min\{\|g\|_\diamond,c\}.
\label{eq:clipped-output-norm}
\end{equation}

Suppose, for contradiction, that
\[
y_-\in\Omega_c^-,
\qquad
y_+\in\Omega_c^+.
\]
Then
\[
\left\|
\mathcal C_c^\diamond(\nabla V(y_-))
\right\|_\diamond
=
\|\nabla V(y_-)\|_\diamond
<
c,
\]
whereas
\[
\left\|
\mathcal C_c^\diamond(\nabla V(y_+))
\right\|_\diamond
=
c.
\]
This contradicts \eqref{eq:clipped-output-equal}.
The opposite mixed itinerary is ruled out identically.
\end{proof}

\begin{remark}[Structural consequence for the post-contact problem]
\label{rem:no-mixed-consequence}
Proposition~\ref{prop:no-mixed-period-two} rules out the
mixed smooth--clipped period-two continuation that one might expect
from a generic switching system.

Thus the post-contact problem for radial gradient clipping is more
constrained: if a period-two branch persists beyond the switching
contact, both points must enter the clipped regime together.

The corresponding post-contact continuation problem will therefore
be formulated as a simultaneous two-point switching bifurcation,
rather than as a mixed-itinerary continuation.
\end{remark}

% ---------------------------------------------------------
\subsection{Summary of the switching geometry}
\label{subsec:switching-geometry-summary}
% ---------------------------------------------------------

The preceding analysis establishes three structural facts that will
be used in the next section.

First, the period-two gradients are exactly opposite:
\begin{equation}
\nabla V(x_+(\mu))
=
-\nabla V(x_-(\mu)).
\label{eq:summary-opposite}
\end{equation}

Second, their common clipping exposure satisfies
\begin{equation}
\mathscr E(\mu)
=
\lambda_*A_*
\|v_*\|_\diamond
\sqrt{\mu}
+
O(\mu^{3/2}).
\label{eq:summary-exposure}
\end{equation}

Third, the first switching interaction is necessarily simultaneous
at both points of the cycle.

Therefore the first-contact condition reduces to the single scalar
equation
\begin{equation}
\mathscr E(\mu_{\mathrm{sc}}(c))
=
c.
\label{eq:first-contact-scalar}
\end{equation}

Section~\ref{sec:separation-law} will invert
\eqref{eq:first-contact-scalar} and derive the resulting separation
law between the smooth flip threshold $\eta_f$ and the first
switching-contact parameter.



% =========================================================
\section{The flip-to-switching separation law}
\label{sec:separation-law}
% =========================================================

We now determine the parameter value at which the smooth
period-two branch created at the flip bifurcation first reaches the
gradient-clipping boundary.

Recall that
\[
\mu=\eta-\eta_f,
\qquad
\eta_f=\frac{2}{\lambda_*},
\]
and let
\[
\mathcal O_2(\mu)
=
\{x_-(\mu),x_+(\mu)\}
\]
denote the smooth period-two branch constructed in
Theorem~\ref{thm:period-two-branch}.

By Proposition~\ref{prop:exact-gradient-balance}, the two points of
the cycle have exactly opposite gradients:
\[
\nabla V(x_+(\mu))
=
-\nabla V(x_-(\mu)).
\]
Hence they have a common clipping exposure
\begin{equation}
\mathscr E(\mu)
:=
\|\nabla V(x_+(\mu))\|_\diamond
=
\|\nabla V(x_-(\mu))\|_\diamond.
\label{eq:section5-exposure}
\end{equation}

The first switching contact therefore occurs when the single scalar
condition
\begin{equation}
\mathscr E(\mu)=c
\label{eq:contact-equation}
\end{equation}
is satisfied.

The main purpose of this section is to solve
\eqref{eq:contact-equation} asymptotically as $c\to0^+$.

Quadratic tangency between period-doubling and border-collision
loci is known for generic piecewise-smooth maps
\cite{SimpsonMeiss2009}.
Accordingly, the exponent two obtained below is not claimed as a
generic new border-collision phenomenon.

The contribution here is the explicit gradient-induced coefficient,
the simultaneous nature of the two-point contact, and the improved
even-power remainder resulting from the exact gradient-balance
structure.

% ---------------------------------------------------------
\subsection{Factorization of the gradient exposure}
\label{subsec:exposure-factorization}
% ---------------------------------------------------------

The square-root law in
Theorem~\ref{thm:exposure-expansion} can be strengthened into a
factorized form that is particularly convenient for inversion.

\begin{lemma}[Exposure factorization]
\label{lem:exposure-factorization}
Assume {\rm (H1)--(H5)}. There exist $\mu_0>0$ and a positive
$C^1$ function
\[
\beta:[0,\mu_0)\to(0,\infty)
\]
such that
\[
E(\mu)=\sqrt{\mu}\,\beta(\mu),
\qquad
0\leq\mu<\mu_0,
\]
with
\[
\beta(0)
=
B_*
=
\lambda_*A_*\|v_*\|_\diamond.
\]
Consequently,
\[
B_*^2
=
\frac{3\lambda_*^4\|v_*\|_\diamond^2}{\Gamma_*}.
\]
\end{lemma}

\begin{proof}
Introduce the signed square-root parameter
\[
\varepsilon^2=\mu.
\]
By Proposition~\ref{prop:refined-cycle-expansion}, the
half-difference of the period-two orbit satisfies
\[
d(-\varepsilon)=-d(\varepsilon)
\]
and
\[
d(\varepsilon)
=
A_*\varepsilon v_*
+
O(\varepsilon^3).
\]

The signed parametrization of
Section~\ref{subsec:square-root-parity} shows that $d$ is a $C^3$
odd function of $\varepsilon$.  Hence it can be written locally as
\begin{equation}
d(\varepsilon)
=
\varepsilon q(\varepsilon^2),
\label{eq:d-factorization}
\end{equation}
where $q$ is a $C^1$ function and
\begin{equation}
q(0)=A_*v_*.
\label{eq:q-zero}
\end{equation}

By the exact gradient-balance identity
\eqref{eq:exact-exposure-d},
\[
\mathscr E(\mu)
=
\frac{2}{\eta_f+\mu}
\|d(\sqrt{\mu})\|_\diamond.
\]
Substituting \eqref{eq:d-factorization} gives
\[
\mathscr E(\mu)
=
\sqrt{\mu}
\frac{2}{\eta_f+\mu}
\|q(\mu)\|_\diamond.
\]
Thus define
\begin{equation}
\beta(\mu)
:=
\frac{2}{\eta_f+\mu}
\|q(\mu)\|_\diamond.
\label{eq:beta-definition}
\end{equation}

Since
\[
q(0)=A_*v_*\neq0
\]
and the norm $\|\cdot\|_\diamond$ is smooth away from the
origin by {\rm (H5)}, the function $\beta$ is $C^1$ after
reducing $\mu_0$ if necessary.

Furthermore,
\begin{align}
\beta(0)
&=
\frac{2}{\eta_f}
\|A_*v_*\|_\diamond
\nonumber\\
&=
\lambda_*A_*\|v_*\|_\diamond
=
B_*,
\end{align}
because
\[
\frac{2}{\eta_f}=\lambda_*.
\]

Finally,
\[
A_*^2
=
\frac{3\lambda_*^2}{\Gamma_*}
\]
by Proposition~\ref{prop:amplitude-identity}. Hence
\[
B_*^2
=
\lambda_*^2A_*^2\|v_*\|_\diamond^2
=
\frac{
3\lambda_*^4\|v_*\|_\diamond^2
}{
\Gamma_*
}.
\]
\end{proof}

\begin{remark}
\label{rem:exposure-factorization-importance}
The factorization
\[
\mathscr E(\mu)
=
\sqrt{\mu}\,\beta(\mu)
\]
is stronger than a formal leading-order expansion.
It separates the universal square-root bifurcation scale from a
regular geometric factor $\beta(\mu)$.

This form will allow the first-contact problem to be reduced to a
regular implicit-function problem in the variable $c^2$.
\end{remark}

% ---------------------------------------------------------
\subsection{Monotonicity and uniqueness of the first contact}
\label{subsec:contact-uniqueness}
% ---------------------------------------------------------

We next show that the exposure increases strictly along the newborn
period-two branch.

\begin{proposition}[Local monotonicity of the exposure]
\label{prop:exposure-monotonicity}
There exists $\mu_1>0$ such that
\begin{equation}
\frac{d}{d\mu}\mathscr E(\mu)>0
\qquad
\text{for all }
0<\mu<\mu_1.
\label{eq:exposure-monotonicity}
\end{equation}
Consequently, $\mathscr E$ is strictly increasing on
$(0,\mu_1)$.
\end{proposition}

\begin{proof}
By Lemma~\ref{lem:exposure-factorization},
\[
\mathscr E(\mu)
=
\sqrt{\mu}\,\beta(\mu),
\]
where
\[
\beta(0)=B_*>0.
\]
For $\mu>0$,
\begin{equation}
\mathscr E'(\mu)
=
\frac{\beta(\mu)}{2\sqrt{\mu}}
+
\sqrt{\mu}\,\beta'(\mu).
\label{eq:exposure-derivative}
\end{equation}

Since $\beta$ is continuous and $\beta(0)=B_*>0$, there exists
$\mu_1>0$ such that
\[
\beta(\mu)\geq\frac{B_*}{2}
\]
for $0\leq\mu<\mu_1$.
Moreover, $\beta'$ is bounded on this interval. Therefore the first
term in \eqref{eq:exposure-derivative}, which is of order
$\mu^{-1/2}$, dominates the second term for sufficiently small
$\mu>0$.

Hence
\[
\mathscr E'(\mu)>0
\]
after reducing $\mu_1$ if necessary.
\end{proof}

We may therefore define the first-contact parameter unambiguously.

\begin{definition}[First switching-contact parameter]
\label{def:first-contact}
For sufficiently small $c>0$, let
$\mu_{\mathrm{sc}}(c)$ denote the unique positive solution of
\begin{equation}
\mathscr E(\mu_{\mathrm{sc}}(c))=c.
\label{eq:mu-sc-definition}
\end{equation}
The corresponding step size is
\begin{equation}
\eta_{\mathrm{sc}}(c)
:=
\eta_f+\mu_{\mathrm{sc}}(c).
\label{eq:eta-sc-definition}
\end{equation}
\end{definition}

% ---------------------------------------------------------
\subsection{The quadratic separation law}
\label{subsec:quadratic-separation}
% ---------------------------------------------------------

We now state the central result of the paper.

\begin{theorem}[Flip-to-switching separation law]
\label{thm:quadratic-separation}
Assume \textbf{(H1)--(H5)}.
There exists $c_0>0$ such that for every
\[
0<c<c_0
\]
the first-contact parameter
$\eta_{\mathrm{sc}}(c)$ is well defined and unique.

As $c\to0^+$,
\begin{equation}
\boxed{
\eta_{\mathrm{sc}}(c)-\eta_f
=
K_\diamond c^2
+
O(c^4)
}
\label{eq:main-separation-law}
\end{equation}
where
\begin{equation}
\boxed{
K_\diamond
=
\frac{1}{
\lambda_*^2A_*^2
\|v_*\|_\diamond^2
}
}
\label{eq:K-A}
\end{equation}
or, equivalently,
\begin{equation}
\boxed{
K_\diamond
=
\frac{\ell_*}{
\lambda_*^3
\|v_*\|_\diamond^2
}
}
\label{eq:K-ell}
\end{equation}
and
\begin{equation}
\boxed{
K_\diamond
=
\frac{\Gamma_*}{
3\lambda_*^4
\|v_*\|_\diamond^2
}.
}
\label{eq:K-Gamma}
\end{equation}

In particular,
\begin{equation}
\lim_{c\to0^+}
\frac{
\eta_{\mathrm{sc}}(c)-\eta_f
}{
c^2
}
=
K_\diamond>0.
\label{eq:K-limit}
\end{equation}
\end{theorem}

\begin{proof}
By Proposition~\ref{prop:exposure-monotonicity}, the exposure
$E(\mu)$ is strictly increasing for all sufficiently small
$\mu>0$. Since
\[
E(\mu)\to0
\qquad\text{as}\qquad
\mu\to0^+,
\]
continuity and strict monotonicity imply that, for every
sufficiently small $c>0$, there exists a unique
$\mu_{\mathrm{sc}}(c)>0$ satisfying
\[
E(\mu_{\mathrm{sc}}(c))=c.
\]

By Theorem~\ref{thm:exposure-expansion},
\[
E(\mu)
=
B_*\sqrt{\mu}
+
O(\mu^{3/2})
=
\sqrt{\mu}\bigl(B_*+O(\mu)\bigr).
\]
Therefore, at $\mu=\mu_{\mathrm{sc}}(c)$,
\[
c^2
=
\mu_{\mathrm{sc}}(c)
\bigl(B_*+O(\mu_{\mathrm{sc}}(c))\bigr)^2,
\]
and hence
\[
c^2
=
B_*^2\mu_{\mathrm{sc}}(c)
+
O\bigl(\mu_{\mathrm{sc}}(c)^2\bigr).
\]

Since
\[
\frac{E(\mu)}{\sqrt{\mu}}
\to B_*>0,
\]
there exists $\mu_0>0$ such that
\[
E(\mu)\geq\frac{B_*}{2}\sqrt{\mu}
\]
for all $0<\mu<\mu_0$.
Evaluating this inequality at
$\mu=\mu_{\mathrm{sc}}(c)$ gives
\[
\mu_{\mathrm{sc}}(c)=O(c^2).
\]
Substituting this estimate into the preceding relation yields
\[
c^2
=
B_*^2\mu_{\mathrm{sc}}(c)
+
O(c^4).
\]
Consequently,
\[
\mu_{\mathrm{sc}}(c)
=
\frac{1}{B_*^2}c^2
+
O(c^4).
\]

Since
\[
\eta_{\mathrm{sc}}(c)
=
\eta_f+\mu_{\mathrm{sc}}(c),
\]
we obtain
\[
\eta_{\mathrm{sc}}(c)-\eta_f
=
K_\diamond c^2+O(c^4),
\]
where
\[
K_\diamond=\frac{1}{B_*^2}.
\]

Using
\[
B_*=\lambda_*A_*\|v_*\|_\diamond
\]
gives
\[
K_\diamond
=
\frac{1}
{\lambda_*^2A_*^2\|v_*\|_\diamond^2}.
\]
Since
\[
A_*^2=\frac{\lambda_*}{\ell_*},
\]
we further obtain
\[
K_\diamond
=
\frac{\ell_*}
{\lambda_*^3\|v_*\|_\diamond^2}.
\]
Finally,
\[
\ell_*=\frac{\Gamma_*}{3\lambda_*},
\]
and therefore
\[
K_\diamond
=
\frac{\Gamma_*}
{3\lambda_*^4\|v_*\|_\diamond^2}.
\]
By {\rm (H4)}, $\Gamma_*>0$, so
\[
K_\diamond>0.
\]
\end{proof}

% ---------------------------------------------------------
\subsection{Explicit tensor form of the separation coefficient}
\label{subsec:explicit-K}
% ---------------------------------------------------------

The coefficient in
Theorem~\ref{thm:quadratic-separation} can be written entirely in
terms of the local differential geometry of the potential.

\begin{corollary}[Explicit separation coefficient]
\label{cor:explicit-K}
Under the assumptions of
Theorem~\ref{thm:quadratic-separation},
\begin{equation}
\boxed{
K_\diamond
=
\frac{
3T_*[v_*,v_*,H_*^{-1}g_*]
-
Q_*[v_*,v_*,v_*,v_*]
}{
3\lambda_*^4
\|v_*\|_\diamond^2
},
}
\label{eq:explicit-K-tensor}
\end{equation}
where
\[
g_*
=
T_*[v_*,v_*,\cdot]^\sharp.
\]

Equivalently, in an orthonormal Hessian eigenbasis
$\{v_i\}_{i=1}^d$,
\begin{align}
K_\diamond
&=
\frac{1}{
3\lambda_*^4
\|v_*\|_\diamond^2
}
\Bigg[
3
\sum_{i=1}^d
\frac{
T_*[v_*,v_*,v_i]^2
}{
\lambda_i
}
\nonumber\\
&\hspace{3em}
-
Q_*[v_*,v_*,v_*,v_*]
\Bigg].
\label{eq:explicit-K-eigenbasis}
\end{align}
\end{corollary}

\begin{proof}
Substitute the representations
\eqref{eq:Gamma-star} and
\eqref{eq:Gamma-eigenbasis}
into \eqref{eq:K-Gamma}.
\end{proof}

\begin{remark}[Geometric interpretation]
\label{rem:K-interpretation}
The coefficient $K_\diamond$ combines three distinct pieces of
local geometry.

First,
\[
\lambda_*
\]
determines the linear flip threshold and the leading displacement of
the gradient in the critical direction.

Second,
\[
\Gamma_*
\]
determines the nonlinear saturation of the unstable critical mode
and hence the amplitude of the emerging period-two orbit.

Third,
\[
\|v_*\|_\diamond
\]
measures the critical eigendirection in the norm used by the
clipping rule.

Thus the switching-contact scale is determined jointly by the
smooth bifurcation geometry and the geometry of the clipping norm.
\end{remark}

% ---------------------------------------------------------
\subsection{Euclidean gradient clipping}
\label{subsec:euclidean-separation}
% ---------------------------------------------------------

For standard gradient clipping,
\[
\|\cdot\|_\diamond=\|\cdot\|_2.
\]
Since
\[
\|v_*\|_2=1,
\]
the separation law simplifies substantially.

\begin{corollary}[Euclidean clipping]
\label{cor:euclidean-separation}
For the standard Euclidean clipping rule,
\begin{equation}
\boxed{
\eta_{\mathrm{sc}}(c)
-
\frac{2}{\lambda_*}
=
K_2c^2
+
O(c^4),
}
\label{eq:euclidean-main-law}
\end{equation}
where
\begin{equation}
\boxed{
K_2
=
\frac{\ell_*}{\lambda_*^3}
=
\frac{\Gamma_*}{3\lambda_*^4}.
}
\label{eq:K-euclidean}
\end{equation}

Equivalently,
\begin{equation}
\boxed{
K_2
=
\frac{1}{3\lambda_*^4}
\left[
3
\sum_{i=1}^d
\frac{
T_*[v_*,v_*,v_i]^2
}{
\lambda_i
}
-
Q_*[v_*,v_*,v_*,v_*]
\right].
}
\label{eq:K-euclidean-expanded}
\end{equation}
\end{corollary}

% ---------------------------------------------------------
\subsection{The same coefficient controls two phenomena}
\label{subsec:coefficient-duality}
% ---------------------------------------------------------

The separation coefficient admits a direct interpretation in terms
of the smooth flip coefficient.

\begin{corollary}[Flip--switching coefficient relation]
\label{cor:flip-switching-relation}
The first flip coefficient $\ell_*$ and the leading
switching-separation coefficient satisfy
\begin{equation}
\boxed{
K_\diamond
=
\frac{\ell_*}{
\lambda_*^3
\|v_*\|_\diamond^2
}.
}
\label{eq:flip-switching-coefficient}
\end{equation}

Therefore, within the present gradient structure, the same nonlinear
coefficient that determines the birth and stability of the smooth
period-two branch also determines its leading-order distance from
the clipping boundary in parameter space.
\end{corollary}

\begin{remark}[Interpretation for Euclidean clipping]
\label{rem:coefficient-duality-euclidean}
For Euclidean clipping,
\[
K_2
=
\frac{\ell_*}{\lambda_*^3}.
\]
Hence a larger positive supercritical flip coefficient produces a
smaller initial oscillation amplitude,
\[
A_*^2=\frac{\lambda_*}{\ell_*},
\]
and therefore delays the first clipping contact to a larger value of
$\eta-\eta_f$.

Conversely, a smaller positive $\ell_*$ produces a larger
small-amplitude cycle for the same $\mu$ and causes clipping to be
encountered sooner.
\end{remark}

% ---------------------------------------------------------
\subsection{Dynamical regimes around the first contact}
\label{subsec:regimes-around-contact}
% ---------------------------------------------------------

The first-contact parameter separates two distinct local regimes.

\begin{proposition}[Pre-contact and contact regimes]
\label{prop:precontact-contact-regimes}
There exist $c_0>0$ and $\mu_0>0$ such that for every
$0<c<c_0$ the following statements hold.

\begin{enumerate}[label=\textnormal{(\roman*)}]

\item
If
\[
0<\mu<\mu_{\mathrm{sc}}(c),
\]
then
\begin{equation}
x_\pm(\mu)\in\Omega_c^-.
\label{eq:precontact-smooth}
\end{equation}
Consequently,
$\mathcal O_2(\mu)$ is simultaneously a stable period-two cycle of
the smooth map $G_{\eta_f+\mu}$ and of the clipped map
$F_{\eta_f+\mu,c}$.

\item
At
\[
\mu=\mu_{\mathrm{sc}}(c),
\]
both points satisfy
\begin{equation}
x_+(\mu_{\mathrm{sc}})
\in\Sigma_c,
\qquad
x_-(\mu_{\mathrm{sc}})
\in\Sigma_c.
\label{eq:both-on-switching}
\end{equation}
The same two points still satisfy the period-two equations for
$F_{\eta_{\mathrm{sc}},c}$.

\item
If
\[
\mu_{\mathrm{sc}}(c)<\mu<\mu_0,
\]
then the continuation of the \emph{smooth reference branch}
satisfies
\begin{equation}
x_\pm(\mu)\in\Omega_c^+.
\label{eq:postcontact-reference-clipped}
\end{equation}
In general, however, this smooth reference cycle is no longer a
period-two cycle of the clipped map.

\end{enumerate}
\end{proposition}

\begin{proof}
By Proposition~\ref{prop:exposure-monotonicity},
the exposure $E(\mu)$ is strictly increasing for all
sufficiently small $\mu>0$. We choose $\mu_0>0$ small enough
that this monotonicity holds throughout $(0,\mu_0)$.

If
\[
0<\mu<\mu_{\mathrm{sc}}(c),
\]
then
\[
E(\mu)<E(\mu_{\mathrm{sc}}(c))=c.
\]
By Theorem~\ref{thm:simultaneous-contact},
both points therefore lie in the smooth region:
\[
x_\pm(\mu)\in\Omega_c^-.
\]
Since $F_{\eta,c}$ and $G_\eta$ coincide at both cycle points,
the smooth period-two cycle is also a period-two cycle of the
clipped map.

At
\[
\mu=\mu_{\mathrm{sc}}(c),
\]
we have
\[
E(\mu_{\mathrm{sc}}(c))=c.
\]
Hence
\[
x_+(\mu_{\mathrm{sc}})\in\Sigma_c,
\qquad
x_-(\mu_{\mathrm{sc}})\in\Sigma_c.
\]
Moreover, the clipping operator satisfies
\[
C_c^\diamond(g)=g
\qquad
\text{whenever }
\|g\|_\diamond=c.
\]
Thus the clipped and smooth maps still agree at both cycle
points, and the period-two relations remain valid at the
contact parameter.

Finally, if
\[
\mu_{\mathrm{sc}}(c)<\mu<\mu_0,
\]
strict monotonicity gives
\[
E(\mu)>c.
\]
Therefore
\[
x_\pm(\mu)\in\Omega_c^+.
\]
These points still form the continuation of the smooth reference
cycle, but in general they no longer form a period-two cycle of
the clipped map, because in $\Omega_c^+$ the gradient update is
rescaled by the factor
\[
\frac{c}{\|\nabla V(x)\|_\diamond}<1.
\]
\end{proof}

\begin{remark}[Stability at the contact point]
\label{rem:contact-stability-not-claimed}
At the exact switching-contact parameter, the two-cycle continues
to exist as an orbit of the clipped map because the map itself is
continuous across $\Sigma_c$.

We do not, however, infer its stability directly from the smooth
multiplier formula. The derivative of the clipped map has different
one-sided limits at the switching surface, so stability at and
beyond contact is a genuinely piecewise-smooth problem.

This issue is deferred to the post-contact analysis.
\end{remark}
\subsection{Origin of the improved remainder}
\label{subsec:origin-improved-remainder}

The $O(c^4)$ remainder in
Theorem~\ref{thm:quadratic-separation}
is a structural consequence of the exact gradient-balance
identity.

Indeed, Theorem~\ref{thm:exposure-expansion} gives
\[
E(\mu)
=
B_*\sqrt{\mu}
+
O(\mu^{3/2}),
\]
rather than the more general expansion
\[
E(\mu)
=
B_1\sqrt{\mu}
+
B_2\mu
+
O(\mu^{3/2}).
\]
The absence of the $O(\mu)$ term is a consequence of the
oddness of the half-difference in the signed square-root
parameter and of the exact identity
\[
E(\mu)
=
\frac{2}{\eta_f+\mu}
\|d(\sqrt{\mu})\|_\diamond .
\]

Squaring the actual exposure expansion gives
\[
E(\mu)^2
=
B_*^2\mu
+
O(\mu^2).
\]
At first contact,
\[
E(\mu_{\mathrm{sc}}(c))=c,
\]
and therefore
\[
c^2
=
B_*^2\mu_{\mathrm{sc}}(c)
+
O\bigl(\mu_{\mathrm{sc}}(c)^2\bigr).
\]
Since
\[
\mu_{\mathrm{sc}}(c)=O(c^2),
\]
it follows that
\[
\mu_{\mathrm{sc}}(c)
=
\frac{1}{B_*^2}c^2
+
O(c^4).
\]

Thus the improvement from the naively expected $O(c^3)$
remainder to $O(c^4)$ does not require a higher-order normal-form
coefficient. It follows directly from the exact two-cycle
gradient balance and the resulting parity structure of the
exposure.

A further expansion of the form
\[
\eta_{\mathrm{sc}}(c)
=
\eta_f
+
K_\diamond c^2
+
K_{\diamond,4}c^4
+
o(c^4)
\]
would require additional higher-order information on the
period-two branch. Such coefficients are not needed for the
leading separation law considered here.
% =========================================================
\section{Post-contact degeneracy and fully clipped period-two dynamics}
\label{sec:post-contact}
% =========================================================

Sections~\ref{sec:switching-geometry} and
\ref{sec:separation-law} show that the smooth period-two branch
created by the supercritical flip reaches the switching hypersurface
simultaneously at both of its points.

The radial clipping structure imposes a further restriction:
a genuine period-two orbit cannot have one point in the smooth region
and the other in the clipped region. Hence any period-two orbit that
persists beyond the first switching contact must enter the clipped
regime at both points.

The post-contact results obtained below are not consequences of the
generic border-collision theory cited in
Section~\ref{sec:introduction}.  They rely specifically on the
radial structure of the clipping operator.

The purpose of this section is to characterize the universal
structure of such fully clipped period-two orbits.  The main result
is that every fully clipped period-two orbit is necessarily
nonhyperbolic: its two-step monodromy matrix has an exact multiplier
equal to one.

This neutral multiplier is not a consequence of the local flip
normal form.  It arises directly from the radial normalization
inherent in gradient clipping.

% ---------------------------------------------------------
\subsection{Fully clipped period-two orbits}
\label{subsec:fully-clipped-two-cycles}
% ---------------------------------------------------------

We first record the exact algebraic structure of a fully clipped
period-two orbit.

\begin{proposition}[Exact geometry of a fully clipped two-cycle]
\label{prop:fully-clipped-geometry}
Let
\[
\mathcal Y_2
=
\{y_-,y_+\}
\]
be a genuine period-two orbit of
\[
F_{\eta,c}(x)
=
x-\eta
\mathcal C_c^\diamond(\nabla V(x)),
\]
and suppose
\[
y_-,y_+\in\Omega_c^+.
\]
Then
\begin{equation}
\boxed{
\|y_+-y_-\|_\diamond
=
\eta c.
}
\label{eq:fully-clipped-step-length}
\end{equation}

Moreover, there exist
\[
u\in\mathbb R^d,
\qquad
\|u\|_\diamond=1,
\]
and numbers
\[
r_+,r_->c
\]
such that
\begin{align}
y_+-y_-
&=
\eta c\,u,
\label{eq:fully-clipped-displacement}
\\
\nabla V(y_+)
&=
r_+u,
\label{eq:fully-clipped-gradient-plus}
\\
\nabla V(y_-)
&=
-r_-u.
\label{eq:fully-clipped-gradient-minus}
\end{align}
\end{proposition}

\begin{proof}
The period-two relations are
\begin{align}
y_-
&=
y_+
-
\eta
\mathcal C_c^\diamond(\nabla V(y_+)),
\label{eq:fully-cycle-plus}
\\
y_+
&=
y_-
-
\eta
\mathcal C_c^\diamond(\nabla V(y_-)).
\label{eq:fully-cycle-minus}
\end{align}
Therefore
\begin{align}
\mathcal C_c^\diamond(\nabla V(y_+))
&=
\frac{y_+-y_-}{\eta},
\label{eq:fully-output-plus}
\\
\mathcal C_c^\diamond(\nabla V(y_-))
&=
-\frac{y_+-y_-}{\eta}.
\label{eq:fully-output-minus}
\end{align}

Since both points lie in the strictly clipped region,
\[
\left\|
\mathcal C_c^\diamond(\nabla V(y_\pm))
\right\|_\diamond
=
c.
\]
Taking norms in \eqref{eq:fully-output-plus} gives
\[
\|y_+-y_-\|_\diamond
=
\eta c.
\]

Define
\[
u
=
\frac{y_+-y_-}{\eta c}.
\]
Then
\[
\|u\|_\diamond=1.
\]

Because radial clipping preserves direction,
\[
\frac{\nabla V(y_+)}
{\|\nabla V(y_+)\|_\diamond}
=
u
\]
and
\[
\frac{\nabla V(y_-)}
{\|\nabla V(y_-)\|_\diamond}
=
-u.
\]
Setting
\[
r_\pm
=
\|\nabla V(y_\pm)\|_\diamond
\]
gives the result.
\end{proof}

\begin{remark}[Amplitude saturation]
\label{rem:amplitude-saturation-main}
Before contact, the period-two amplitude is generated by the smooth
flip mechanism and satisfies
\[
\|x_+(\mu)-x_-(\mu)\|
=
O(\sqrt{\mu}).
\]
After both points enter the clipped regime, the step length is no
longer determined by the smooth normal form.  Instead, radial clipping
imposes the exact constraint
\[
\|y_+-y_-\|_\diamond
=
\eta c.
\]

Thus clipping replaces the smooth square-root amplitude law by an
exact saturated step-length law.
\end{remark}

% ---------------------------------------------------------
\subsection{Exclusion of mixed period-two itineraries}
\label{subsec:no-mixed-two-cycle-main}
% ---------------------------------------------------------

We now state explicitly the structural restriction established in
Section~\ref{sec:switching-geometry}.

\begin{proposition}[No mixed smooth--clipped period-two orbit]
\label{prop:no-mixed-two-cycle-main}
A radial clipped gradient map cannot possess a genuine period-two
orbit
\[
\{y_-,y_+\}
\]
such that exactly one point lies in $\Omega_c^-$ and the other lies
in $\Omega_c^+$.

Consequently, every genuine period-two orbit is locally of one of the
following types:
\begin{enumerate}[label=\textnormal{(\roman*)}]
\item both points are smooth;
\item both points lie on the switching boundary;
\item both points are strictly clipped.
\end{enumerate}
\end{proposition}

\begin{proof}
Suppose, for contradiction, that
\[
y_-\in\Omega_c^-,
\qquad
y_+\in\Omega_c^+.
\]
Then the period-two equations imply
\[
\nabla V(y_-)
+
\mathcal C_c^\diamond(\nabla V(y_+))
=
0.
\]
Taking the clipping norm gives
\[
\|\nabla V(y_-)\|_\diamond
=
\left\|
\mathcal C_c^\diamond(\nabla V(y_+))
\right\|_\diamond
=
c,
\]
which contradicts
\[
\|\nabla V(y_-)\|_\diamond<c.
\]
The opposite mixed itinerary is ruled out identically.
\end{proof}

\begin{remark}
\label{rem:no-mixed-global}
Proposition~\ref{prop:no-mixed-two-cycle-main} is not a local
consequence of the flip bifurcation.  It follows from the algebraic
structure of radial clipping and therefore applies to arbitrary
period-two orbits, independently of their amplitude or origin.
\end{remark}

% ---------------------------------------------------------
\subsection{Differential structure in the clipped region}
\label{subsec:clipped-differential-main}
% ---------------------------------------------------------
We now examine the differential structure of the radial clipping
operator in the strictly clipped region. This structure will be
the key ingredient in the spectral analysis of fully clipped
period-two orbits.

Let
\[
\rho(g):=\|g\|_\diamond .
\]
For any nonzero vector $g\in\mathbb{R}^d$, write
\[
u:=\frac{g}{\rho(g)},
\qquad
\rho(u)=1,
\]
and define the covector
\[
\phi_u:=D\rho(u).
\]

Since $\rho$ is positively homogeneous of degree one,
Euler's identity gives
\begin{equation}
\phi_u[u]
=
D\rho(u)[u]
=
\rho(u)
=
1.
\label{eq:main-euler-rho}
\end{equation}
Moreover, the derivative of a positively homogeneous function
of degree one is homogeneous of degree zero. Hence, since
$g=\rho(g)u$ with $\rho(g)>0$,
\begin{equation}
D\rho(g)=D\rho(u)=\phi_u.
\label{eq:main-rho-derivative-homogeneity}
\end{equation}

In the strictly clipped region
\[
\rho(g)>c,
\]
the clipping operator takes the form
\[
C_c^\diamond(g)
=
c\frac{g}{\rho(g)}.
\]
For any $h\in\mathbb{R}^d$, differentiation yields
\[
DC_c^\diamond(g)[h]
=
\frac{c}{\rho(g)}h
-
\frac{c}{\rho(g)^2}
g\,D\rho(g)[h].
\]
Using
\[
g=\rho(g)u
\]
and \eqref{eq:main-rho-derivative-homogeneity}, we obtain
\[
DC_c^\diamond(g)[h]
=
\frac{c}{\rho(g)}
\left(
h-u\,\phi_u[h]
\right).
\]
Equivalently,
\begin{equation}
DC_c^\diamond(g)
=
\frac{c}{\rho(g)}
\left(
I-u\otimes\phi_u
\right).
\label{eq:main-clipping-derivative}
\end{equation}

The operator
\[
P_u:=I-u\otimes\phi_u
\]
has a simple geometric interpretation. By
\eqref{eq:main-euler-rho},
\[
P_u u
=
u-u\phi_u[u]
=
0,
\]
and
\[
\phi_u P_u
=
\phi_u-\phi_u[u]\phi_u
=
0.
\]
Thus $P_u$ projects onto
\[
T_u\{\rho=1\}
=
\ker\phi_u,
\]
the tangent space of the unit $\rho$-sphere at $u$, along the
radial direction spanned by $u$.

In particular, multiplying
\eqref{eq:main-clipping-derivative} from the left by $\phi_u$
gives
\begin{equation}
\phi_uDC_c^\diamond(g)
=
\frac{c}{\rho(g)}
\left(
\phi_u-\phi_u[u]\phi_u
\right)
=
0.
\label{eq:main-annihilation}
\end{equation}

Equation \eqref{eq:main-annihilation} expresses the differential
degeneracy created by radial clipping: infinitesimal variations
in the gradient-magnitude direction are annihilated, whereas only
variations tangent to the clipping sphere contribute to the
differential of the clipped update.

For every $x\in\Omega_c^+$, the chain rule therefore gives
\begin{equation}
DF_{\eta,c}(x)
=
I
-
\eta
DC_c^\diamond(\nabla V(x))
D^2V(x).
\label{eq:main-clipped-Jacobian}
\end{equation}

\begin{proposition}[Pointwise neutral covector]
\label{prop:pointwise-neutral-covector}
Let $x\in\Omega_c^+$ and set
\[
u
=
\frac{\nabla V(x)}
{\|\nabla V(x)\|_\diamond}.
\]
Then
\begin{equation}
\phi_uDF_{\eta,c}(x)=\phi_u.
\label{eq:main-pointwise-neutral}
\end{equation}
In particular,
\[
1\in\sigma\bigl(DF_{\eta,c}(x)\bigr).
\]
\end{proposition}

\begin{proof}
Using \eqref{eq:main-clipped-Jacobian},
\[
\phi_uDF_{\eta,c}(x)
=
\phi_u
-
\eta
\phi_u
DC_c^\diamond(\nabla V(x))
D^2V(x).
\]
By \eqref{eq:main-annihilation},
\[
\phi_u
DC_c^\diamond(\nabla V(x))
=
0.
\]
Hence the second term vanishes and
\[
\phi_uDF_{\eta,c}(x)=\phi_u.
\]
Thus $\phi_u$ is a left eigenvector of $DF_{\eta,c}(x)$
associated with the eigenvalue $1$. Therefore
\[
1\in\sigma\bigl(DF_{\eta,c}(x)\bigr).
\]
\end{proof}

\begin{remark}[Pointwise neutrality versus orbit neutrality]
\label{rem:pointwise-versus-orbit-neutrality}
Although every one-step Jacobian in the strictly clipped region
has a unit eigenvalue, the corresponding left eigenvector depends
on the local gradient direction. Hence, along a general clipped
orbit, unit eigenvalues of the individual Jacobians do not imply
that their product possesses a unit Floquet multiplier.

Period-two orbits are exceptional. As shown in
Proposition~7.1, the clipped period-two equations force the two
gradient directions to be exactly opposite. Moreover, because
$\rho$ is even,
\[
\rho(-u)=\rho(u),
\]
and therefore
\[
D\rho(-u)=-D\rho(u).
\]
Consequently,
\[
(-u)\otimes D\rho(-u)
=
u\otimes D\rho(u),
\]
so the tangent projection
\[
I-u\otimes\phi_u
\]
is unchanged under $u\mapsto -u$.

Thus the same neutral covector structure is preserved at both
points of a fully clipped period-two orbit. This observation is
the mechanism behind the universal unit-multiplier theorem proved
in the next subsection.
\end{remark}
                                                                                         


% ---------------------------------------------------------
\subsection{Universal unit multiplier for fully clipped two-cycles}
\label{subsec:universal-unit-main}
% ---------------------------------------------------------

We now obtain the principal post-contact result.

\begin{theorem}[Universal unit multiplier]

\label{thm:universal-unit-main}
Let
\[
\mathcal Y_2
=
\{y_-,y_+\}
\]
be any fully clipped period-two orbit of $F_{\eta,c}$.

Let
\[
M_c
:=
DF_{\eta,c}(y_-)
DF_{\eta,c}(y_+)
\label{eq:main-monodromy}
\]
be the corresponding two-step monodromy matrix.

Then
\begin{equation}
\boxed{
1\in\sigma(M_c).
}
\label{eq:main-unit-multiplier}
\end{equation}

More precisely, if $u$ is the common clipping direction defined by
Proposition~\ref{prop:fully-clipped-geometry}, then
\begin{equation}
\varphi_u
DF_{\eta,c}(y_+)
=
\varphi_u
\label{eq:common-left-plus}
\end{equation}
and
\begin{equation}
\varphi_u
DF_{\eta,c}(y_-)
=
\varphi_u.
\label{eq:common-left-minus}
\end{equation}
Hence
\begin{equation}
\varphi_uM_c
=
\varphi_u.
\label{eq:common-left-monodromy}
\end{equation}
\end{theorem}


\begin{proof}
By Proposition~\ref{prop:fully-clipped-geometry}, there exist
$r_+,r_->c$ and a vector $u$ with $\|u\|_\diamond=1$ such that
\[
\nabla V(y_+)=r_+u,
\qquad
\nabla V(y_-)=-r_-u.
\]

Since every norm is even,
\[
\rho(-z)=\rho(z),
\]
and therefore
\[
D\rho(-u)=-D\rho(u).
\]
Thus
\[
\phi_{-u}=-\phi_u.
\]

At $y_+$, the normalized gradient direction is $u$.
Hence Proposition~\ref{prop:pointwise-neutral-covector} gives
\[
\phi_u DF_{\eta,c}(y_+)=\phi_u.
\]

At $y_-$, the normalized gradient direction is $-u$.
Applying Proposition~\ref{prop:pointwise-neutral-covector}
with $-u$ gives
\[
\phi_{-u}DF_{\eta,c}(y_-)=\phi_{-u}.
\]
Since $\phi_{-u}=-\phi_u$, this is equivalent to
\[
\phi_uDF_{\eta,c}(y_-)=\phi_u.
\]

Consequently,
\[
\phi_uM_c
=
\phi_uDF_{\eta,c}(y_-)
DF_{\eta,c}(y_+)
=
\phi_u.
\]
Thus $\phi_u$ is a nonzero left eigenvector of $M_c$
associated with the eigenvalue $1$. Therefore
\[
1\in\sigma(M_c).
\]
\end{proof}

\begin{corollary}[Fully clipped period-two cycles are nonhyperbolic]
\label{cor:fully-clipped-nonhyperbolic}
Every fully clipped period-two orbit of a radial clipped gradient map
is nonhyperbolic.
\end{corollary}

\begin{proof}
Theorem~\ref{thm:universal-unit-main} gives a multiplier equal to
$1$, which lies on the unit circle.
\end{proof}

\begin{corollary}[No linear exponential attraction]
\label{cor:no-linear-exponential-attraction}
A fully clipped period-two orbit cannot be linearly exponentially
asymptotically stable.
\end{corollary}

\begin{proof}
Linear exponential asymptotic stability requires every Floquet
multiplier of the period-two orbit to lie strictly inside the unit
disk.  Theorem~\ref{thm:universal-unit-main} provides an exact
multiplier equal to $1$.
\end{proof}

\begin{remark}[Nonhyperbolic does not mean unstable]
\label{rem:nonhyperbolic-not-unstable}
The presence of the unit multiplier does not by itself imply
nonlinear instability.

A fully clipped period-two orbit may be nonlinearly attracting,
repelling, neutral, or embedded in a family of periodic orbits.
These alternatives depend on higher-order dynamics in the neutral
direction and cannot be distinguished by the linearization alone.
\end{remark}

% ---------------------------------------------------------
\subsection{The spectral transition at first contact}
\label{subsec:spectral-transition}
% ---------------------------------------------------------

Before contact, the smooth period-two orbit is hyperbolically
attracting along its critical direction.

By Theorem~\ref{thm:period-two-stability},
\begin{equation}
\rho_{\mathrm{sm}}^c(\mu)
=
1
-
4\lambda_*\mu
+
O(\mu^{3/2}).
\label{eq:precontact-critical-multiplier}
\end{equation}

At the first switching contact,
\begin{equation}
\mu_{\mathrm{sc}}(c)
=
K_\diamond c^2
+
O(c^4).
\label{eq:contact-mu-main}
\end{equation}
Hence
\begin{equation}
\rho_{\mathrm{sm}}^c
\bigl(
\mu_{\mathrm{sc}}(c)
\bigr)
=
1
-
4\lambda_*K_\diamond c^2
+
O(c^3).
\label{eq:contact-smooth-multiplier}
\end{equation}

Since
\[
K_\diamond
=
\frac{
\Gamma_*
}{
3\lambda_*^4
\|v_*\|_\diamond^2
},
\]
this becomes
\begin{equation}
\rho_{\mathrm{sm}}^c
\bigl(
\mu_{\mathrm{sc}}(c)
\bigr)
=
1
-
\frac{
4\Gamma_*
}{
3\lambda_*^3
\|v_*\|_\diamond^2
}
c^2
+
O(c^3).
\label{eq:contact-smooth-multiplier-Gamma}
\end{equation}

Thus the smooth cycle reaches the switching boundary while still
weakly attracting:
\[
\rho_{\mathrm{sm}}^c<1.
\]

Any fully clipped period-two orbit beyond contact, however, satisfies
\[
\rho_{\mathrm{clip}}^c=1
\]
in at least one direction.

This yields the following conclusion.

\begin{corollary}[Loss of hyperbolicity under fully clipped continuation]
\label{cor:loss-hyperbolicity-clipping}

Suppose that a period-two orbit admits a continuation from the
first switching contact into the fully clipped regime.
Then every such fully clipped period-two continuation is
necessarily nonhyperbolic.

Immediately before contact, the smooth critical multiplier satisfies
\[
\rho_{\mathrm{sm}}^c
=
1-4\lambda_*K_\diamond c^2+O(c^3),
\]
and therefore
\[
1-\rho_{\mathrm{sm}}^c
=
4\lambda_*K_\diamond c^2+O(c^3).
\]
Thus the smooth period-two branch reaches the switching boundary
while still weakly attracting, whereas, by
Theorem~\ref{thm:universal-unit-main}, every fully clipped
period-two orbit carries an exact unit multiplier.
\end{corollary}

\begin{remark}
\label{rem:spectral-transition-not-small-jump}
The transition should not be interpreted as a regular perturbation
of the smooth two-cycle multiplier.

The derivative of radial clipping changes discontinuously across
the switching boundary.  On the clipped side, radial variations of
the gradient are annihilated, and this rank-deficient structure
produces an exact neutral direction.

Thus the post-contact spectrum is constrained by the geometry of
the clipping operator rather than obtained by smoothly perturbing
the pre-contact smooth monodromy.
\end{remark}

% ---------------------------------------------------------
\subsection{Degeneracy of the period-two continuation equation}
\label{subsec:continuation-degeneracy}
% ---------------------------------------------------------

A period-two point of the clipped map is a solution of
\begin{equation}
\mathcal P(y,\eta,c)
=
F_{\eta,c}^2(y)-y
=
0.
\label{eq:period-two-continuation-equation}
\end{equation}

If the two-cycle were hyperbolic, one would normally attempt to
continue it by applying the implicit function theorem to
\[
\mathcal P=0.
\]
The universal unit multiplier prevents this.

\begin{proposition}[Singularity of the post-contact continuation problem]
\label{prop:continuation-singularity}
Let $y$ belong to a fully clipped period-two orbit. Then
\begin{equation}
D_y\mathcal P(y,\eta,c)
=
DF_{\eta,c}^2(y)-I
\label{eq:period-two-equation-Jacobian}
\end{equation}
is singular.

Consequently, assumptions \textbf{(H1)--(H5)} do not yield a
standard implicit-function continuation of the period-two branch
through the fully clipped regime.
\end{proposition}

\begin{proof}
By Theorem~\ref{thm:universal-unit-main},
\[
1
\in
\sigma
\left(
DF_{\eta,c}^2(y)
\right).
\]
Hence
\[
0
\in
\sigma
\left(
DF_{\eta,c}^2(y)-I
\right),
\]
so the derivative of the period-two fixed-point equation is
singular.
\end{proof}

\begin{remark}[Why a universal post-contact branch theorem is absent]
\label{rem:no-universal-postcontact-theorem}
The smooth period-two branch before contact is governed by a
nondegenerate flip bifurcation and therefore admits a standard local
branch description.

After clipping becomes active, the relevant periodic-orbit equation
is intrinsically degenerate.  Additional nonlinear information is
required to determine whether the fully clipped periodic orbit:

\begin{itemize}
\item persists as an isolated nonhyperbolic two-cycle;
\item belongs to a family of two-cycles;
\item is nonlinearly attracting or repelling;
\item or terminates at the switching contact.
\end{itemize}

The solvable planar model of Section~\ref{sec:planar-model} realizes
the second possibility: the post-contact dynamics contain a
one-parameter family of fully clipped period-two cycles.
\end{remark}

% ---------------------------------------------------------
\subsection{Euclidean clipping}
\label{subsec:euclidean-postcontact-main}
% ---------------------------------------------------------

For the Euclidean norm,
\[
\|\cdot\|_\diamond
=
\|\cdot\|_2,
\]
the geometry becomes especially transparent.

If
\[
u
=
\frac{\nabla V(x)}
{\|\nabla V(x)\|_2},
\]
then in the clipped region
\begin{equation}
D\mathcal C_c(\nabla V(x))
=
\frac{c}{\|\nabla V(x)\|_2}
\left(
I-uu^\top
\right).
\label{eq:euclidean-clipping-Jacobian-main}
\end{equation}

Hence
\begin{equation}
DF_{\eta,c}(x)
=
I
-
\frac{\eta c}{
\|\nabla V(x)\|_2
}
\left(
I-uu^\top
\right)
D^2V(x).
\label{eq:euclidean-clipped-map-Jacobian}
\end{equation}

The radial covector satisfies
\begin{equation}
u^\top DF_{\eta,c}(x)
=
u^\top.
\label{eq:euclidean-radial-left-eigenvector}
\end{equation}

For a fully clipped period-two orbit with gradient directions
$\pm u$,
\begin{equation}
u^\top
DF_{\eta,c}(y_-)
DF_{\eta,c}(y_+)
=
u^\top.
\label{eq:euclidean-two-step-neutral}
\end{equation}

Thus the unit multiplier has a direct geometric interpretation:
Euclidean clipping removes the derivative of the update in the
radial gradient-magnitude direction.

% ---------------------------------------------------------
\subsection{Summary of the post-contact mechanism}
\label{subsec:postcontact-summary-main}
% ---------------------------------------------------------

The local transition established in this paper can therefore be
summarized as follows.

For
\[
0<
\eta-\eta_f
<
\eta_{\mathrm{sc}}(c)-\eta_f,
\]
the period-two orbit lies entirely in the smooth region and is
governed by the supercritical smooth flip.

At
\[
\eta
=
\eta_{\mathrm{sc}}(c),
\]
both points of the orbit reach the switching boundary simultaneously.

A mixed smooth--clipped two-cycle cannot occur.

Any genuine period-two continuation into the clipped region must
satisfy
\[
\|y_+-y_-\|_\diamond
=
\eta c
\]
and necessarily possesses an exact unit multiplier:
\[
1
\in
\sigma(M_c).
\]

Hence the first clipping contact marks a transition
\begin{equation}
\begin{aligned}
\text{smooth attracting two-cycle}
&\\[-1mm]
\Downarrow
&\\[-1mm]
\text{simultaneous clipping contact}
&\\[-1mm]
\Downarrow
&\\[-1mm]
\text{nonhyperbolic clipped regime}.
\end{aligned}
\end{equation}
This degeneracy is a structural consequence of radial clipping and
does not depend on the specific solvable model considered later.
% =========================================================
\section{A completely solvable planar gradient model}
\label{sec:planar-model}
% =========================================================

We now present a two-dimensional model for which the smooth flip
bifurcation, the first clipping contact, the separation coefficient,
and the post-contact dynamics can all be computed explicitly.

Consider the potential
\begin{equation}
V(x,y)
=
\frac{\lambda}{2}x^2
+
\frac{\nu}{2}y^2
-
\frac{b}{4}x^4
+
\frac{d}{6}x^6,
\label{eq:planar-potential}
\end{equation}
where
\begin{equation}
\lambda>\nu>0,
\qquad
b>0,
\qquad
d>0,
\qquad
b^2<4d\lambda.
\label{eq:planar-parameters}
\end{equation}

We use the standard Euclidean clipping norm throughout this section.

The gradient is
\begin{equation}
\nabla V(x,y)
=
\begin{pmatrix}
\lambda x-bx^3+dx^5
\\
\nu y
\end{pmatrix}.
\label{eq:planar-gradient}
\end{equation}
For convenience, define
\begin{equation}
g(x)
:=
\lambda x-bx^3+dx^5
=
x\bigl(\lambda-bx^2+dx^4\bigr).
\label{eq:g-definition}
\end{equation}

The condition $b^2<4d\lambda$ implies
\begin{equation}
\lambda-bs+ds^2>0
\qquad
\text{for every }s\geq0,
\label{eq:g-positive-factor}
\end{equation}
and hence
\begin{equation}
xg(x)>0
\qquad
\text{for every }x\neq0.
\label{eq:g-sign}
\end{equation}

Thus $(0,0)$ is the unique critical point of $V$ and is a strict
global minimum.

% ---------------------------------------------------------
\subsection{Verification of the standing hypotheses}
\label{subsec:planar-hypotheses}
% ---------------------------------------------------------

At the origin,
\begin{equation}
H_*
=
D^2V(0,0)
=
\begin{pmatrix}
\lambda&0
\\
0&\nu
\end{pmatrix}.
\label{eq:planar-Hessian}
\end{equation}
Since
\[
\lambda>\nu>0,
\]
the maximal Hessian eigenvalue is
\begin{equation}
\lambda_*=\lambda
\label{eq:planar-lambda-star}
\end{equation}
and its normalized eigenvector is
\begin{equation}
v_*
=
e_1
=
\begin{pmatrix}
1\\0
\end{pmatrix}.
\label{eq:planar-v-star}
\end{equation}

The flip threshold is therefore
\begin{equation}
\boxed{
\eta_f
=
\frac{2}{\lambda}.
}
\label{eq:planar-flip-threshold}
\end{equation}

Since the potential is even in $x$,
\begin{equation}
D^3V(0,0)=0.
\label{eq:planar-third-zero}
\end{equation}
Moreover,
\begin{equation}
D^4V(0,0)[e_1,e_1,e_1,e_1]
=
-6b.
\label{eq:planar-fourth}
\end{equation}

Hence
\begin{equation}
\Gamma_*
=
-
D^4V(0,0)[e_1^4]
=
6b>0.
\label{eq:planar-Gamma}
\end{equation}
Thus the flip is supercritical.

The first flip coefficient is
\begin{equation}
\boxed{
\ell_*
=
\frac{\Gamma_*}{3\lambda}
=
\frac{2b}{\lambda}.
}
\label{eq:planar-ell}
\end{equation}

The general amplitude formula gives
\begin{equation}
\boxed{
A_*^2
=
\frac{3\lambda^2}{\Gamma_*}
=
\frac{\lambda^2}{2b},
}
\label{eq:planar-A-square}
\end{equation}
or
\begin{equation}
A_*
=
\frac{\lambda}{\sqrt{2b}}.
\label{eq:planar-A}
\end{equation}

Finally, the general Euclidean separation coefficient becomes
\begin{equation}
\boxed{
K_2
=
\frac{\Gamma_*}{3\lambda^4}
=
\frac{2b}{\lambda^4}.
}
\label{eq:planar-K-general-prediction}
\end{equation}

We shall recover all of these quantities below from the exact
dynamics.

% ---------------------------------------------------------
\subsection{Exact smooth gradient dynamics}
\label{subsec:planar-smooth-dynamics}
% ---------------------------------------------------------

The ordinary gradient-descent map is
\begin{equation}
G_\eta(x,y)
=
\begin{pmatrix}
(1-\eta\lambda)x
+
\eta bx^3
-
\eta dx^5
\\[1mm]
(1-\eta\nu)y
\end{pmatrix}.
\label{eq:planar-GD-map}
\end{equation}

The line
\begin{equation}
\mathcal L
=
\{(x,0):x\in\mathbb R\}
\label{eq:invariant-axis}
\end{equation}
is invariant.

On $\mathcal L$, the dynamics reduce to the odd scalar map
\begin{equation}
f_\eta(x)
=
(1-\eta\lambda)x
+
\eta bx^3
-
\eta dx^5.
\label{eq:scalar-smooth-map}
\end{equation}

Because $f_\eta$ is odd, a symmetric period-two orbit has the form
\[
\{-a,a\},
\qquad
a>0.
\]

The condition
\[
f_\eta(a)=-a
\]
is equivalent to
\begin{equation}
\lambda-\frac{2}{\eta}
=
ba^2-da^4.
\label{eq:smooth-cycle-equation}
\end{equation}

Set
\begin{equation}
q:=a^2.
\label{eq:q-amplitude}
\end{equation}
Then
\begin{equation}
dq^2-bq
+
\lambda-\frac{2}{\eta}
=
0.
\label{eq:q-quadratic}
\end{equation}

The small-amplitude root is therefore
\begin{equation}
q_{\mathrm{sm}}(\eta)
=
\frac{
b-
\sqrt{
b^2
-
4d\left(\lambda-\frac{2}{\eta}\right)
}
}{
2d}.
\label{eq:q-smooth-exact}
\end{equation}

Thus the smooth period-two branch is given explicitly by
\begin{equation}
\boxed{
\mathcal O_{2,\mathrm{sm}}(\eta)
=
\left\{
(-a_{\mathrm{sm}}(\eta),0),
(a_{\mathrm{sm}}(\eta),0)
\right\},
}
\label{eq:smooth-cycle-exact}
\end{equation}
where
\begin{equation}
a_{\mathrm{sm}}(\eta)
=
\sqrt{q_{\mathrm{sm}}(\eta)}.
\label{eq:a-smooth}
\end{equation}

% ---------------------------------------------------------
\subsection{Exact verification of the square-root law}
\label{subsec:planar-square-root}
% ---------------------------------------------------------

Let
\begin{equation}
\mu
=
\eta-\frac{2}{\lambda}.
\label{eq:planar-mu}
\end{equation}
Then
\begin{equation}
\lambda-\frac{2}{\eta}
=
\frac{\lambda^2}{2}\mu
-
\frac{\lambda^3}{4}\mu^2
+
O(\mu^3).
\label{eq:lambda-two-over-eta-expansion}
\end{equation}

Solving \eqref{eq:smooth-cycle-equation} gives
\begin{equation}
a_{\mathrm{sm}}(\eta)^2
=
\frac{\lambda^2}{2b}\mu
+
\frac{
\lambda^3(d\lambda-b^2)
}{
4b^3
}
\mu^2
+
O(\mu^3).
\label{eq:a-square-expansion}
\end{equation}

Taking the square root yields
\begin{equation}
a_{\mathrm{sm}}(\eta)
=
\frac{\lambda}{\sqrt{2b}}
\mu^{1/2}
\left[
1
+
\frac{
\lambda(d\lambda-b^2)
}{
4b^2
}
\mu
+
O(\mu^2)
\right].
\label{eq:a-expansion}
\end{equation}

Hence
\begin{equation}
a_{\mathrm{sm}}(\eta)
=
A_*\sqrt{\mu}
+
O(\mu^{3/2}),
\label{eq:planar-square-root-law}
\end{equation}
with
\[
A_*=\frac{\lambda}{\sqrt{2b}},
\]
exactly as predicted by
\eqref{eq:planar-A}.

Because the model is reflection-symmetric, the midpoint of the
smooth two-cycle remains exactly at the critical point:
\begin{equation}
\frac{
(a_{\mathrm{sm}},0)
+
(-a_{\mathrm{sm}},0)
}{2}
=
(0,0).
\label{eq:planar-midpoint-zero}
\end{equation}

% ---------------------------------------------------------
\subsection{Stability of the smooth period-two branch}
\label{subsec:planar-smooth-stability}
% ---------------------------------------------------------

The derivative of the scalar map is
\begin{equation}
f_\eta'(x)
=
1-\eta\lambda
+
3\eta bx^2
-
5\eta dx^4.
\label{eq:scalar-derivative}
\end{equation}

Using the cycle equation
\eqref{eq:smooth-cycle-equation}, one obtains
\begin{equation}
f_\eta'(a)
=
-1
+
2\eta a^2
\left(
b-2da^2
\right).
\label{eq:scalar-derivative-cycle}
\end{equation}

Since $f_\eta'$ is even, the critical multiplier of the two-step
map is
\begin{equation}
\rho_x^{\mathrm{sm}}
=
\left[
-1
+
2\eta a^2
\left(
b-2da^2
\right)
\right]^2.
\label{eq:planar-smooth-x-multiplier}
\end{equation}

For $\mu\to0^+$,
\begin{equation}
\rho_x^{\mathrm{sm}}
=
1-4\lambda\mu
+
O(\mu^2),
\label{eq:planar-smooth-x-expansion}
\end{equation}
in agreement with
Theorem~\ref{thm:period-two-stability}.

The transverse multiplier is
\begin{equation}
\rho_y^{\mathrm{sm}}
=
(1-\eta\nu)^2.
\label{eq:planar-smooth-y-multiplier}
\end{equation}

Since $0<\nu<\lambda$, at $\eta=\eta_f$ one has
$|1-2\nu/\lambda|<1$, so the transverse multiplier stays inside the
unit disk.  Consequently, the smooth period-two cycle is locally
asymptotically stable for all sufficiently small $\eta-\eta_f>0$.

% ---------------------------------------------------------
\subsection{Exact gradient exposure}
\label{subsec:planar-exposure}
% ---------------------------------------------------------

At the positive point of the cycle,
\[
\nabla V(a_{\mathrm{sm}},0)
=
\begin{pmatrix}
g(a_{\mathrm{sm}})
\\
0
\end{pmatrix}.
\]
By the exact period-two relation,
\begin{equation}
g(a_{\mathrm{sm}})
=
\frac{2a_{\mathrm{sm}}}{\eta}.
\label{eq:planar-exact-gradient-balance}
\end{equation}

Hence the common Euclidean gradient exposure is
\begin{equation}
\boxed{
\mathscr E(\eta)
=
\frac{2a_{\mathrm{sm}}(\eta)}{\eta}.
}
\label{eq:planar-exposure-exact}
\end{equation}

Using \eqref{eq:a-expansion},
\begin{equation}
\mathscr E(\mu)
=
\frac{\lambda^2}{\sqrt{2b}}
\mu^{1/2}
\left[
1
+
\frac{
\lambda(d\lambda-3b^2)
}{
4b^2
}
\mu
+
O(\mu^2)
\right].
\label{eq:planar-exposure-expansion}
\end{equation}

In particular,
\begin{equation}
B_*
=
\frac{\lambda^2}{\sqrt{2b}},
\label{eq:planar-B}
\end{equation}
and therefore
\begin{equation}
\frac{1}{B_*^2}
=
\frac{2b}{\lambda^4},
\label{eq:planar-B-inverse}
\end{equation}
again agreeing with the general separation coefficient
\eqref{eq:planar-K-general-prediction}.

% ---------------------------------------------------------
\subsection{Exact switching-contact equation}
\label{subsec:planar-contact-equation}
% ---------------------------------------------------------

The first contact occurs when
\begin{equation}
\frac{2a_{\mathrm{sm}}(\eta)}{\eta}
=
c.
\label{eq:planar-contact-condition}
\end{equation}
Hence at contact
\begin{equation}
a_{\mathrm{sm}}
=
\frac{\eta c}{2}.
\label{eq:planar-contact-amplitude}
\end{equation}

Substituting this into
\eqref{eq:smooth-cycle-equation} gives the exact scalar equation
\begin{equation}
\Psi(\eta,c)
:=
\lambda
-
\frac{2}{\eta}
-
\frac{b}{4}\eta^2c^2
+
\frac{d}{16}\eta^4c^4
=
0.
\label{eq:Psi-contact}
\end{equation}

At
\[
(\eta,c)
=
\left(
\frac{2}{\lambda},0
\right),
\]
we have
\[
\Psi\left(\frac{2}{\lambda},0\right)=0
\]
and
\begin{equation}
\partial_\eta\Psi
\left(
\frac{2}{\lambda},0
\right)
=
\frac{\lambda^2}{2}
>0.
\label{eq:Psi-eta-positive}
\end{equation}

Thus the implicit function theorem gives a unique smooth
first-contact curve
\[
\eta=\eta_{\mathrm{sc}}(c)
\]
for sufficiently small $c$.

% ---------------------------------------------------------
\subsection{Exact verification of the separation law}
\label{subsec:planar-separation}
% ---------------------------------------------------------

Write
\begin{equation}
\eta_{\mathrm{sc}}(c)
=
\frac{2}{\lambda}
+
K_2c^2
+
K_4c^4
+
O(c^6).
\label{eq:planar-eta-series}
\end{equation}

Substituting this expansion into
\eqref{eq:Psi-contact} and comparing powers of $c$ gives
\begin{equation}
\boxed{
K_2
=
\frac{2b}{\lambda^4},
}
\label{eq:planar-K2}
\end{equation}
and
\begin{equation}
\boxed{
K_4
=
\frac{6b^2}{\lambda^7}
-
\frac{2d}{\lambda^6}.
}
\label{eq:planar-K4}
\end{equation}

Therefore
\begin{equation}
\boxed{
\eta_{\mathrm{sc}}(c)
=
\frac{2}{\lambda}
+
\frac{2b}{\lambda^4}c^2
+
\left(
\frac{6b^2}{\lambda^7}
-
\frac{2d}{\lambda^6}
\right)c^4
+
O(c^6).
}
\label{eq:planar-separation-final}
\end{equation}

The leading coefficient
\[
\frac{2b}{\lambda^4}
\]
coincides exactly with
\[
\frac{\Gamma_*}{3\lambda_*^4}
=
\frac{6b}{3\lambda^4},
\]
as predicted by
Theorem~\ref{thm:quadratic-separation}.

This model therefore verifies not only the quadratic exponent but
also the absence of a cubic correction.

% ---------------------------------------------------------
\subsection{The fully clipped symmetric branch}
\label{subsec:planar-clipped-symmetric}
% ---------------------------------------------------------

We now turn to the post-contact regime.  Define
\begin{equation}
a_c(\eta)
:=
\frac{\eta c}{2},
\label{eq:ac-definition}
\end{equation}
and consider the pair
\begin{equation}
Y_\pm(\eta,c)
=
(\pm a_c(\eta),0).
\label{eq:symmetric-clipped-pair}
\end{equation}
Its gradient magnitude is
\begin{equation}
g(a_c)
=
a_c
\left(
\lambda
-
ba_c^2
+
da_c^4
\right),
\label{eq:g-ac}
\end{equation}
and a direct calculation gives
\begin{equation}
g(a_c)-c
=
\frac{\eta c}{2}
\Psi(\eta,c),
\label{eq:g-minus-c-Psi}
\end{equation}
with $\Psi$ as in \eqref{eq:Psi-contact}.  Hence $g(a_c)=c$ at
$\eta=\eta_{\mathrm{sc}}(c)$, and since
$\partial_\eta\Psi(\eta_{\mathrm{sc}}(c),c)>0$ for all sufficiently
small $c$,
\begin{equation}
g(a_c)>c
\qquad
\text{for }
\eta>\eta_{\mathrm{sc}}(c)
\label{eq:ac-fully-clipped}
\end{equation}
sufficiently close to the contact value.  Both points $Y_\pm$
therefore lie in the fully clipped region.  On the invariant axis,
radial clipping acts as
\[
F_{\eta,c}(x,0)
=
(x-\eta c,0)\quad(g(x)>c),
\qquad
F_{\eta,c}(x,0)
=
(x+\eta c,0)\quad(|g(x)|>c,\ x<0),
\]
and since $2a_c=\eta c$,
\begin{equation}
F_{\eta,c}(a_c,0)
=
(-a_c,0),
\qquad
F_{\eta,c}(-a_c,0)
=
(a_c,0).
\label{eq:clipped-symmetric-forward}
\end{equation}

Thus:

\begin{proposition}[Explicit fully clipped period-two branch]
\label{prop:explicit-clipped-branch}
For every sufficiently small $c>0$ and every
$\eta>\eta_{\mathrm{sc}}(c)$ sufficiently close to
$\eta_{\mathrm{sc}}(c)$,
\begin{equation}
\boxed{
\mathcal O_{2,\mathrm{clip}}(\eta,c)
=
\left\{
\left(-\frac{\eta c}{2},0\right),
\left(\frac{\eta c}{2},0\right)
\right\}
}
\label{eq:explicit-clipped-cycle}
\end{equation}
is a fully clipped period-two orbit.

Its amplitude satisfies exactly
\begin{equation}
\boxed{
2a_c=\eta c.
}
\label{eq:exact-amplitude-saturation-planar}
\end{equation}
\end{proposition}

% ---------------------------------------------------------
\subsection{A continuum of clipped period-two cycles}
\label{subsec:continuum-clipped-cycles}
% ---------------------------------------------------------

The symmetric clipped cycle is not isolated.

For
\[
\eta>\eta_{\mathrm{sc}}(c),
\]
we have
\[
g(a_c)>c.
\]
By continuity, there exists $\delta>0$ such that
\begin{equation}
g(a_c+s)>c,
\qquad
g(a_c-s)>c
\label{eq:clipped-neighborhood-condition}
\end{equation}
for every
\[
|s|<\delta.
\]

Define
\begin{align}
Y_+(s)
&=
(a_c+s,0),
\label{eq:Y-plus-s}
\\
Y_-(s)
&=
(-a_c+s,0).
\label{eq:Y-minus-s}
\end{align}

After reducing $\delta$ if necessary,
\[
a_c+s>0,
\qquad
-a_c+s<0.
\]

Since
\[
Y_+(s)-Y_-(s)
=
(\eta c,0),
\]
radial clipping gives
\begin{align}
F_{\eta,c}(Y_+(s))
&=
Y_-(s),
\label{eq:family-forward}
\\
F_{\eta,c}(Y_-(s))
&=
Y_+(s).
\label{eq:family-backward}
\end{align}

We therefore obtain the following stronger result.

\begin{theorem}[A local family of clipped period-two cycles]
\label{thm:family-clipped-cycles}
For every sufficiently small $c>0$ and
$\eta>\eta_{\mathrm{sc}}(c)$ sufficiently close to contact, there
exists $\delta>0$ such that
\begin{equation}
\mathcal O_2(s)
=
\left\{
(-a_c+s,0),
(a_c+s,0)
\right\},
\qquad
|s|<\delta,
\label{eq:period-two-family}
\end{equation}
is a one-parameter family of fully clipped period-two cycles.

Equivalently,
\begin{equation}
F_{\eta,c}^2(x,0)
=
(x,0)
\label{eq:F2-identity-axis}
\end{equation}
on a nontrivial interval of the positive clipped component and on
its image under $F_{\eta,c}$.
\end{theorem}

\begin{proof}
For $|s|<\delta$,
both points lie in the clipped regime by
\eqref{eq:clipped-neighborhood-condition}.

At $Y_+(s)$, the gradient points in the positive $x$ direction, so
the clipped step is exactly
\[
-\eta c e_1.
\]
Therefore
\[
F_{\eta,c}(Y_+(s))
=
(a_c+s-\eta c,0)
=
(-a_c+s,0)
=
Y_-(s).
\]

At $Y_-(s)$, the gradient points in the negative $x$ direction, so
the clipped step is
\[
+\eta c e_1.
\]
Thus
\[
F_{\eta,c}(Y_-(s))
=
(-a_c+s+\eta c,0)
=
(a_c+s,0)
=
Y_+(s).
\]
\end{proof}

\begin{remark}[Geometric meaning of the unit multiplier]
\label{rem:family-unit-meaning}
The continuum in
Theorem~\ref{thm:family-clipped-cycles} gives a concrete nonlinear
interpretation of the universal unit multiplier proved in
Theorem~\ref{thm:universal-unit-main}.

The neutral multiplier is not merely an artifact of linearization:
in this model the two-step map is exactly the identity along the
one-dimensional clipped family.
\end{remark}

% ---------------------------------------------------------
\subsection{Exact post-contact multipliers}
\label{subsec:planar-post-multipliers}
% ---------------------------------------------------------

Consider the period-two cycle
\[
\mathcal O_2(s)
=
\{Y_-(s),Y_+(s)\}.
\]
Set
\begin{equation}
r_+(s)
=
g(a_c+s),
\qquad
r_-(s)
=
g(a_c-s).
\label{eq:r-plus-minus}
\end{equation}

At either point, the Euclidean clipping direction is parallel to
$e_1$. Therefore
\begin{equation}
I-e_1e_1^\top
=
\begin{pmatrix}
0&0\\
0&1
\end{pmatrix}.
\label{eq:e1-projector}
\end{equation}

Since
\[
D^2V(x,0)
=
\begin{pmatrix}
\lambda-3bx^2+5dx^4&0
\\
0&\nu
\end{pmatrix},
\]
the one-step Jacobians along the clipped cycle are
\begin{align}
DF_{\eta,c}(Y_+(s))
&=
\begin{pmatrix}
1&0
\\
0&
1-\dfrac{\eta\nu c}{r_+(s)}
\end{pmatrix},
\label{eq:Jac-plus-family}
\\
DF_{\eta,c}(Y_-(s))
&=
\begin{pmatrix}
1&0
\\
0&
1-\dfrac{\eta\nu c}{r_-(s)}
\end{pmatrix}.
\label{eq:Jac-minus-family}
\end{align}

Hence the two-step monodromy matrix is
\begin{equation}
M_{\mathrm{clip}}(s)
=
\begin{pmatrix}
1&0
\\
0&
\displaystyle
\left(
1-\frac{\eta\nu c}{r_+(s)}
\right)
\left(
1-\frac{\eta\nu c}{r_-(s)}
\right)
\end{pmatrix}.
\label{eq:family-monodromy}
\end{equation}

Thus the two multipliers are
\begin{equation}
\boxed{
\rho_x^{\mathrm{clip}}=1
}
\label{eq:planar-unit-multiplier}
\end{equation}
and
\begin{equation}
\boxed{
\rho_y^{\mathrm{clip}}(s)
=
\left(
1-\frac{\eta\nu c}{r_+(s)}
\right)
\left(
1-\frac{\eta\nu c}{r_-(s)}
\right).
}
\label{eq:planar-transverse-multiplier}
\end{equation}

For the symmetric member $s=0$,
\begin{equation}
\rho_y^{\mathrm{clip}}(0)
=
\left(
1-\frac{\eta\nu c}{g(a_c)}
\right)^2.
\label{eq:symmetric-clipped-transverse}
\end{equation}

At the switching contact,
\[
g(a_c)=c,
\]
so the outer limiting transverse multiplier becomes
\begin{equation}
\rho_y^{\mathrm{clip}}
=
(1-\eta_{\mathrm{sc}}\nu)^2.
\label{eq:transverse-contact-limit}
\end{equation}

Since
\[
\eta_{\mathrm{sc}}(c)
\to
\frac{2}{\lambda}
\]
as $c\to0^+$ and
\[
0<\nu<\lambda,
\]
one has
\[
\left|
1-\eta_{\mathrm{sc}}(c)\nu
\right|
<1
\]
for all sufficiently small $c$.

Hence the clipped family is transversely attracting while remaining
exactly neutral along its one-dimensional period-two direction.

% ---------------------------------------------------------
\subsection{The local transition in the solvable model}
\label{subsec:planar-transition-summary}
% ---------------------------------------------------------

The planar model provides an explicit realization of the entire
smooth-to-clipped scenario.  For $\eta<\eta_f=2/\lambda$ the origin
is locally asymptotically stable; for
$\eta_f<\eta<\eta_{\mathrm{sc}}(c)$ there is a unique small stable
smooth two-cycle
$\{(-a_{\mathrm{sm}},0),(a_{\mathrm{sm}},0)\}$ whose amplitude
grows as
\[
a_{\mathrm{sm}}
\sim
\frac{\lambda}{\sqrt{2b}}
\sqrt{\eta-\eta_f}.
\]
At $\eta=\eta_{\mathrm{sc}}(c)$ both points reach the clipping
boundary simultaneously with
\[
a_{\mathrm{sm}}
=
\frac{\eta_{\mathrm{sc}}c}{2},
\qquad
\eta_{\mathrm{sc}}(c)
=
\frac{2}{\lambda}
+
\frac{2b}{\lambda^4}c^2
+
O(c^4).
\label{eq:planar-summary-separation}
\]
For $\eta>\eta_{\mathrm{sc}}(c)$ sufficiently close to contact, the
clipped map possesses a one-parameter family of period-two cycles
with exact step constraint $x_+-x_-=\eta c$, and the smooth critical
multiplier $\rho_x^{\mathrm{sm}}<1$ is replaced by the exact neutral
multiplier $\rho_x^{\mathrm{clip}}=1$.  The model thus exhibits the
transition from a stable smooth period-two branch, through a
simultaneous clipping contact, to a neutral family of clipped
two-cycles.

\subsection{A concrete parameter choice}

To make the preceding formulas explicit, consider the parameter
choice
\[
\lambda=2,\qquad
\nu=\frac12,\qquad
b=1,\qquad
d=1.
\]
Then
\[
\eta_f=1,
\qquad
\Gamma_*=6,
\qquad
\ell_*=1,
\]
and therefore
\[
A_*=\sqrt{2},
\qquad
B_*=2\sqrt{2},
\qquad
K_2=\frac18.
\]
Hence the smooth period-two branch satisfies
$x_\pm(\mu)=x_*\pm\sqrt{2\mu}\,v_*+O(\mu)$, while the first
switching threshold obeys
\[
\eta_{\mathrm{sc}}(c)
=
1+\frac18c^2+O(c^4),
\]
with the next correction computable explicitly:
\[
\eta_{\mathrm{sc}}(c)
=
1+\frac18c^2+\frac1{64}c^4+O(c^6).
\]
The example thus provides an explicit realization of the general
quadratic separation law of Section~\ref{sec:separation-law}.
% =========================================================
\section{Conclusion and discussion}
\label{sec:conclusion}
% =========================================================

We studied the first interaction between a smooth flip bifurcation
and the endogenous switching geometry generated by radial gradient
clipping.  The mechanism operates on two sharply separated scales:
clipping is locally invisible near the critical point and becomes
dynamically relevant only at finite amplitude.

\paragraph{Local scale.}
For every fixed $c>0$ the clipped map $F_{\eta,c}$ coincides with the
ordinary gradient map $G_\eta$ on a neighborhood of a nondegenerate
critical point; hence the fixed point, its linearization, its local
invariant structures, its center-manifold reduction, and its smooth
flip coefficient are independent of the clipping threshold.  In
particular, for a strict local minimum $x_*$ with simple maximal
Hessian eigenvalue $\lambda_*$, the first stability threshold remains
\begin{equation}
\eta_f
=
\frac{2}{\lambda_*}.
\label{eq:conclusion-eta-f}
\end{equation}
Under the supercriticality condition
\begin{equation}
\Gamma_*
=
3T_*[v_*,v_*,H_*^{-1}g_*]
-
Q_*[v_*^4]
>0,
\label{eq:conclusion-Gamma}
\end{equation}
the loss of stability creates a stable smooth period-two branch
\begin{equation}
x_\pm(\mu)
=
x_*
\pm
A_*\sqrt{\mu}\,v_*
+
O(\mu),
\qquad
\mu=\eta-\eta_f,
\qquad
A_*^2
=
\frac{3\lambda_*^2}{\Gamma_*}.
\label{eq:conclusion-amplitude}
\end{equation}

\paragraph{Finite-amplitude scale.}
Clipping becomes active only when this orbit reaches the switching
hypersurface
\[
\Sigma_c
=
\{x:\|\nabla V(x)\|_\diamond=c\}.
\]
The gradient structure makes the interaction rigid: every smooth
period-two orbit satisfies the exact identity
\begin{equation}
\nabla V(x_+)
=
-\nabla V(x_-),
\label{eq:conclusion-gradient-balance}
\end{equation}
so both points have identical clipping exposure and reach $\Sigma_c$
simultaneously.  Matching the square-root amplitude scale with the
threshold produces the principal asymptotic law
\begin{equation}
\boxed{
\eta_{\mathrm{sc}}(c)-\eta_f
=
\frac{\Gamma_*}{
3\lambda_*^4
\|v_*\|_\diamond^2
}
c^2
+
O(c^4),
}
\label{eq:conclusion-main-law}
\end{equation}
which for Euclidean clipping becomes
\begin{equation}
\boxed{
\eta_{\mathrm{sc}}(c)
=
\frac{2}{\lambda_*}
+
\frac{\Gamma_*}{
3\lambda_*^4
}
c^2
+
O(c^4).
}
\label{eq:conclusion-euclidean}
\end{equation}
The even-power expansion reflects an additional symmetry of the
gradient problem: in the signed square-root parameter
$\varepsilon=\sqrt{\eta-\eta_f}$ the half-difference of the two-cycle
is odd while its midpoint is even, which together with
\eqref{eq:conclusion-gradient-balance} removes the cubic correction
from the contact curve.

\paragraph{Post-contact scale.}
A clipped period-two orbit cannot mix smooth and clipped points, so
any continuation past the first contact enters the clipped regime at
both points simultaneously, and every fully clipped period-two orbit
obeys the exact step-length constraint
\begin{equation}
\|y_+-y_-\|_\diamond
=
\eta c.
\label{eq:conclusion-step-law}
\end{equation}
Most importantly, radial clipping creates an unavoidable neutral
direction: the monodromy matrix
\[
M_c
=
DF_{\eta,c}(y_-)
DF_{\eta,c}(y_+)
\]
of a fully clipped period-two orbit satisfies
\begin{equation}
\boxed{
1\in\sigma(M_c).
}
\label{eq:conclusion-unit}
\end{equation}
Every fully clipped two-cycle is therefore nonhyperbolic: the radial
normalization removes first-order information in the
gradient-magnitude direction.  The solvable planar model of
Section~\ref{sec:planar-model} shows that this unit multiplier is not
a formal artifact --- it corresponds to an actual loss of isolation of
the periodic orbit.

% ---------------------------------------------------------
\subsection{Dynamical interpretation}
\label{subsec:conclusion-interpretation}
% ---------------------------------------------------------

The results separate two effects that are often conflated.  The
first is the loss of local stability of the minimum: a smooth
gradient-descent phenomenon governed by $D^2V(x_*)$, $D^3V(x_*)$,
$D^4V(x_*)$, and unaffected by clipping.  The second is the
modification of an already finite-amplitude trajectory once it
reaches $\|\nabla V\|_\diamond=c$.  Clipping does not postpone the
first local flip threshold; it introduces a second dynamical scale
beyond it, the switching-contact scale
$\eta_{\mathrm{sc}}(c)-\eta_f=O(c^2)$.  This distinction is useful
when interpreting large-step optimization dynamics: linear stability
determines when the critical point loses stability, while the
clipping threshold determines when the resulting nonlinear oscillation
begins to experience a different piecewise-smooth evolution.

% ---------------------------------------------------------
\subsection{Limitations}
\label{subsec:conclusion-limitations}
% ---------------------------------------------------------

The theory focuses on an isolated nondegenerate minimum with a simple
maximal Hessian eigenvalue, which isolates a one-dimensional flip mode
but excludes the additional neutral directions that arise when the set
of minimizers forms a manifold.  Second, the post-contact unit
multiplier prevents a universal hyperbolic continuation theorem:
under the present assumptions alone, the existence, uniqueness, and
nonlinear stability of a fully clipped continuation are not
determined.  Third, the paper treats radial clipping; coordinatewise
clipping, adaptive normalization, momentum, and stochastic mini-batch
updates generate different switching geometries.  Finally, the
analysis is local in the bifurcation parameter and does not classify
the remote global dynamics after the orbit enters the clipped region.

% ---------------------------------------------------------
\subsection{Directions for further work}
\label{subsec:conclusion-future}
% ---------------------------------------------------------

Several extensions appear natural: reducing the nonhyperbolic
post-contact dynamics to a higher-order scalar or center-manifold
normal form, which the exact unit multiplier suggests may form a new
class of degenerate border interactions specific to radial gradient
maps; replacing the isolated minimum by a Morse--Bott manifold of
minima, combining the $-1$ flip direction with $+1$ tangent
directions; degenerate smooth flips with $\Gamma_*=0$, where a
matching principle $r(\mu)\asymp\mu^\gamma$, $d(c)\asymp c^p$
suggests the general contact scale
$\mu_{\mathrm{sc}}\asymp c^{p/\gamma}$; and determining which
structures persist for momentum, stochastic gradients, and nonradial
clipping.

% ---------------------------------------------------------
\subsection{Final perspective}
\label{subsec:conclusion-final-perspective}
% ---------------------------------------------------------

Near a nondegenerate minimum, radial gradient clipping is locally
invisible: $F_{\eta,c}=G_\eta$, so the smooth map determines when the
first oscillatory instability is born.  Once the oscillation reaches
finite amplitude, the same potential generates an endogenous switching
boundary at the explicit scale
\[
\eta_{\mathrm{sc}}(c)-\eta_f
\sim
\frac{\Gamma_*}{
3\lambda_*^4
\|v_*\|_\diamond^2
}
c^2.
\]
At that point the dynamics cease to be governed solely by the smooth
flip normal form: radial clipping removes the gradient-magnitude
degree of freedom and forces the post-contact period-two dynamics to
carry an exact neutral multiplier.  The clipped gradient map thus
exhibits a sharp transition from smooth bifurcation theory to
intrinsically nonhyperbolic switching dynamics, with the transition
scale determined explicitly by the local geometry of the underlying
potential.


\section*{Acknowledgements}
During the preparation of this work, the authors used the DeepSeek
V4 Harness (deepseek-v4-flash model) for \LaTeX{} typesetting,
compilation, and manuscript-preparation assistance, and OpenAI
ChatGPT (GPT-5.6 Sol) for mathematical discussion, language editing,
and manuscript organization.  All mathematical arguments,
references, and conclusions were independently checked by the
authors, who take full responsibility for the manuscript.

\section*{Data availability statement}
Data sharing is not applicable to this article as no datasets were
generated or analysed during the current study.

\section*{Conflict of interest}
The authors declare that they have no conflict of interest.

\section*{Author contributions}
Both authors contributed to the conception, analysis, and writing of
the manuscript.

\section*{Funding}
The authors declare that no specific funding was received for this
work.

\bibliographystyle{iopart-num}
\bibliography{references}

\end{document}
